% Grafo publico acumulativo de apendices HMT--MD.
% Exclusiones vinculantes: este archivo no incorpora el radion ni las tablas
% abreviadas de diez o doce masas. Tampoco alcanza A5_suplemento_digital_rev11.
% Todos los \input de este archivo se verifican como archivos existentes.

\chapter{Certificados algebraicos y combinatorios de los objetos finitos}
\label{app:infraestructura-reproducible-acumulativa}
\label{app:pruebas}

Los objetos finitos utilizados en las demostraciones admiten un segundo canal
de comprobación. Este canal reconstruye catálogos, operadores e invariantes
desde las semillas declaradas y los contrasta con identidades exactas o cotas
dirigidas; complementa las demostraciones, pero no las sustituye.

\HMTDemoteHeadings
\chapter{Catálogos finitos y certificados de reproducibilidad}
\label{app:datos-reproducibles}

Los objetos finitos siguientes permiten repetir las verificaciones sin
incorporar notación de producción a la exposición matemática. Cada control
conserva el dominio, la operación, el invariante y el falsador que le
corresponden.

\section{Catálogos finitos del estado común}

\begin{longtable}{@{}>{\raggedright\arraybackslash}p{.34\textwidth}
  >{\raggedright\arraybackslash}p{.58\textwidth}@{}}
\toprule
Objeto finito & Contenido matemático verificable\\
\midrule
\endhead
Catálogo de emisiones de longitud seis
& Las \(468\) emisiones \(U_6\), sus multiplicidades, orientación y firmas
aditiva y multiplicativa.\\
Partición espectral de las emisiones
& Las veintiséis clases de tamaño \(18\), junto con el mapa explícito que
reconstruye cada emisión desde su clase y su posición interna.\\
Atlas dodecafásico de rutas
& Las \(468\) rutas mínimas con hoja, fase, acarreo, frontera y memoria.\\
Catálogo del lector arquimediano
& Las cinco regiones de \(\pi\), las regiones conjugadas de \(e\) y el eje de
autoescala de \(\varphi\), con sus cilindros y multiplicidades.\\
Flujo de Hensel de veinte bloques
& La matriz unimodular, su inversa entera, los estados \(3\)-ádicos y la
recurrencia de residuo--cociente que prolonga los tres canales.\\
Registro de monodromía nonádica
& Los veinte bloques, sus fases, prefijos ternarios y prefijos decimales;
incluye los \(43\) retornos completos y los \(387\) pares \((K,K+9)\) por
canal.\\
Ventana de supervivencia \(K=5,\ldots,13\)
& Los candidatos locales, la rama compatible y las tres salidas de la
transición \(9\to1\), sin identificar retorno de fase con reinicio del estado.\\
Lectura arquimediana inicial
& Las doce ternas decimales obtenidas de los primeros trece bloques de cada
canal, separadas del certificado de precisión arbitraria.\\
\bottomrule
\end{longtable}

\section{Identidades algebraicas, operatoriales y espectrales de los certificados finitos}

\begin{longtable}{@{}>{\raggedright\arraybackslash}p{.34\textwidth}
  >{\raggedright\arraybackslash}p{.58\textwidth}@{}}
\toprule
Control & Identidades y falsadores\\
\midrule
\endhead
Monodromía y generación de constantes
& Verifica \(w_6\to w_{30}\to R_{36}\to G_9\), la selección \(3\to1\),
el cambio de estado tras nueve fases, las cinco regiones de \(\pi\) y la
compatibilidad de truncamientos. Falla si se introduce una cifra objetivo
antes del lector.\\
Coordenada angular conjugada
& Comprueba el selector regional \(169\), la recurrencia determinantal, la
suma analítica y su cota de resto, y la unicidad de \(C^*\). Falla si una
constante del SI interviene en la selección.\\
Dinámica del cociente \(468=18\cdot26\)
& Certifica la biyección entre emisiones y clases, el retorno uniforme, la
medida de la sombra promediada, los cociclos de modo y la elevación
microscópica. Distingue la frecuencia cronológica de la medida de Parry.\\
Canales enteros de las constantes
& Verifica las identidades \((729,271,315,161)\), las realizaciones
conjuntistas de \(271/315\), las órbitas diedrales y la factorización de
paridad visible y paridad de acarreo.\\
Incidencia de Witt
& Reconstruye la incidencia punto--hexada, el módulo \(E_{30}\), el
transporte positivo de \(P_3\), la estrella de Witt y las obstrucciones de
equivarianza.\\
Selector variacional
& Reconstruye el normalizador \(S_3\), el promedio de Reynolds, la
factorización polar y la cadena de proyectores de rango tres; la positividad
se prueba exactamente en \(\mathbb Q(\sqrt5)\).\\
Elevación \(W_{12}\to W_{24}\)
& Comprueba las incidencias de cuatro y cinco puntos, la elevación única de
hexadas y la correspondencia tipada con las cinco regiones de \(\pi\).\\
Realizaciones de \(-1/12\)
& Construye los proyectores racionales de los desplazamientos de pasos tres
y cuatro, verifica simetría e idempotencia y obtiene
\((3-4)/12=-1/12\) como diferencia de trazas normalizadas.\\
Holonomía de Schläfli
& Recorre las plaquetas, certifica \(A^2=5I-4A+5J\), la configuración de las
veintisiete rectas, sus \(45\) triángulos y la partición
\(729=27+270+432\).\\
Funcionales espectrales
& Recalcula Catalán, Stieltjes, Apéry, Bendersky--Glaisher--Kinkelin, las
dos extensiones Euler--Kronecker y las cotas dirigidas de productos primos.\\
Vacancias y criticidad
& Comprueba las separaciones \(21/22\), los retornos \(6/7\), la parte
armónica finita y los aproximantes dodecafásicos del modo crítico.\\
\bottomrule
\end{longtable}

\section{Diccionario tipado de los catálogos}

\begin{longtable}{>{\ttfamily}p{.18\textwidth}>{\raggedright\arraybackslash}p{.72\textwidth}}
\toprule
Campo & Significado y codominio\\
\midrule
\endhead
Esig & Firma del canal multiplicativo; identifica una de las veintiséis
clases de la factorización del emisor.\\
hplus\_type & Clase de orientación del canal aditivo para los dos pares de
direcciones cardinales declarados.\\
diag\_c & Índice entero de la clase diagonal del cursor aditivo.\\
px\_size & Cardinal de la fibra de la clase aditiva correspondiente.\\
colw & Sexteto de pesos de columna del bloque emitido; pertenece a
\(\mathbb Z_{\geq0}^{6}\).\\
\bottomrule
\end{longtable}

\section{Reproducibilidad de los certificados y obstrucciones exactas de las construcciones finitas}

Los controles emplean aritmética entera o racional siempre que el objeto lo
permite. Las cotas analíticas usan intervalos con redondeo dirigido y separan
la demostración simbólica de la comparación decimal. Cada ejecución
reconstruye los objetos desde las semillas declaradas, verifica las identidades
de incidencia y naturalidad, y contrasta las salidas con una implementación
aritmética independiente.

El control del producto nativo rechaza cualquier dependencia de una evaluación
decimal, una factorización externa, una criba, una tabla de ceros o una
constante objetivo. El control nonádico rechaza la identificación de la fibra
ambiente de \(729\) elementos con el cardinal de un espacio de historias. El
control de fase rechaza \(g(K+9)=g(K)\) como prueba de igualdad de estados. El
control metrológico rechaza toda selección de ruta realizada después de
consultar una magnitud observada.

Una discrepancia entre una identidad del texto, el catálogo finito y su control
independiente refuta la certificación correspondiente. Esta regla no sustituye
ninguna demostración: fija un segundo canal de comprobación para los objetos
finitos que intervienen en ella.

\HMTRestoreHeadings

\section{Compatibilidad global entre tipos, dominios, morfismos y certificados}
\label{sec:ejecucion-integral-acumulativa}

El control integral comprueba la cadena causal, repite las identidades
matemáticas sobre los objetos finitos y contrasta las realizaciones exactas con
sus formulaciones independientes. Una divergencia entre una definición, su
tabla matemática y el invariante declarado refuta el control correspondiente.

% Pruebas auxiliares y trazabilidad del cuaderno de 64 páginas. La matriz
% humana se conserva como testigo subordinado a la matriz preventiva común;
% la tabla de fuentes se publica como procedencia del aporte, no como índice
% rival del registro canónico vigente.

\chapter{Dependencias exactas entre potencia genealógica, información reversible y obstrucción Gödel–Turing}
\label{app:integracion-vn64-controles}

\HMTDemoteHeadings
\chapter{Dependencias lógicas y certificados de consistencia}

\section{Cadena causal de los certificados: estados finitos \(\to\) restricciones \(\to\) supervivencia \(\to\) sección infinita \(\to\) realizaciones}

La cadena lógica de la integración no parte de una conclusión arquimediana.
Se organiza en el orden siguiente:
\[
\begin{tikzpicture}[baseline=(current bounding box.center),>=Latex,
 node distance=7mm and 8mm,
 n/.style={rounded corners,draw=hmtblue,fill=hmtlightblue,
 minimum height=8mm,align=center,font=\small}]
 \node[n] (finite) {estados finitos\\tipados};
 \node[n,right=of finite] (maps) {restricciones\\compatibles};
 \node[n,right=of maps] (surv) {supervivencia\\y poda};
 \node[n,right=of surv] (lim) {sección\\infinita};
 \node[n,right=of lim] (read) {lectores\\y realizaciones};
 \draw[->,thick] (finite)--(maps);
 \draw[->,thick] (maps)--(surv);
 \draw[->,thick] (surv)--(lim);
 \draw[->,thick] (lim)--(read);
\end{tikzpicture}
\]
Un certificado de cifras verifica una restricción de la última flecha; no
crea retrospectivamente los estados ni sustituye el teorema de existencia
del límite.

\section{Identidades aritméticas exactas de los certificados finitos}

Los censos que acompañan al certificado de filtración nonádica satisfacen
\[
 396\cdot6=2376,
 \qquad
 43\cdot9=387,
 \qquad
 432+4\cdot144=1008.
\]
La primera igualdad relaciona eventos elementales de frontera y trits
emitidos por canal; la segunda cuenta las parejas de fase \((K,K+9)\); la
tercera reconstruye la multiplicidad total de las cinco regiones de
\(\pi\).

Los dos papeles regionales que no pueden permutarse son
\[
\begin{array}{rcl}
U_\pi^\perp&=&(501,614,498,169,272,272),\\
E_{\rm sig}&=&555555,\qquad \text{multiplicidad }432,\\[.35em]
U_\pi^{--}&=&(870,923,810,418,674,521),\\
E_{\rm sig}&=&933717,\qquad \text{multiplicidad }144.
\end{array}
\]

\section{Existencia y unicidad de la rama infinita en un árbol finitamente ramificado de prefijos supervivientes}

\begin{lemma}[Árbol finitamente ramificado]
Sea \(T\) un árbol con niveles finitos no vacíos y tal que todo nodo
superviviente tiene un descendiente superviviente. Entonces existe una rama
infinita. Si, además, cada nodo de la subtorre seleccionada tiene exactamente
un descendiente compatible, la rama es única.
\end{lemma}

\begin{proof}
La existencia es el lema de König aplicado al árbol finitamente ramificado.
Para la unicidad, dos ramas distintas tendrían un primer nivel en el que
divergen, dando dos descendientes compatibles del mismo prefijo, contra la
hipótesis.
\end{proof}

\begin{controlbox}
Una ejecución a profundidad \(N\) sólo prueba lo ocurrido en el árbol
truncado. La afirmación infinita requiere la ley uniforme que preserva no
vaciedad y compatibilidad para todo nivel. El cuerpo separa expresamente
ambos estatutos y, para la filtración nonádica, conserva su selector uniforme
certificado.
\end{controlbox}

\section{Clausura y potencia plena}

Para \(A\subseteq X=\varprojlim X_n\), se ha probado
\[
 \overline A=\bigcap_n p_n^{-1}(p_n(A)).
\]
El conjunto de sucesiones binarias finalmente nulas es denso y propio en
\(2^{\mathbb N}\), pero posee todas las sombras finitas. Por ello cualquier
algoritmo que reciba sólo \((p_n(A))_n\) no puede distinguir \(A\) de su
clausura. Este ejemplo bloquea la falsa identidad
\[
 \varprojlim\mathcal P(X_n)=\mathcal P(\varprojlim X_n).
\]

\section{Compatibilidad tipada de los operadores}

La diferencia entre la inclusión unital y la esquina marcada se verifica
directamente:
\[
\begin{aligned}
 \tau_{d+1}(a\otimes I_9)
 &=\frac{\operatorname{Tr}(a)\operatorname{Tr}(I_9)}{9^{d+1}}
 =\tau_d(a),\\
 \tau_{d+1}(a\otimes p_j)
 &=\frac{\operatorname{Tr}(a)}{9^{d+1}}
 =\frac1{9}\tau_d(a).
\end{aligned}
\]
La primera rama es unital y tracial; la segunda conserva una esquina marcada
y no puede usarse para inferir el factor \(II_1\).

\section{Anulación del coste de Landauer en la inversión reversible del estado}

Para una variable completa \(X\) y una proyección \(Y=q(X)\),
\[
 H(X)=H(Y)+H(X\mid Y).
\]
Si \(q\) es biyectiva, la fibra es unitaria y \(H(X\mid Y)=0\). Si \(q\)
es un reinicio del TRIT uniforme, la información descartada es \(\log3\) y
el coste mínimo es \(k_B\Theta\log3\). Para la extensión
\(\widehat E_N\) a \(L^1\) de una esperanza condicional normal
traza-preservante \(E_N\), y bajo las hipótesis de finitud declaradas,
\[
 S_\tau(\widehat E_Nh)-S_\tau(h)
 =D_\tau(h\Vert \widehat E_Nh)\ge0.
\]
Las dos fórmulas localizan el mismo fenómeno en dominios diferentes: el
coste pertenece a la contracción de una fibra, no a la evolución reversible
del estado completo.

\section{Obstrucción Gödel–Turing a un decisor uniforme de multiplicidades proyectivas}

Sea \(T_e\) la máquina de índice \(e\). En el nivel \(n\) se define
\[
 X_n^e=\{b_n\}\cup
 \{c_n:T_e\text{ no se detuvo antes de }n\}.
\]
Cada \(X_n^e\) es computable. La fibra global tiene dos ramas si \(T_e\) no
se detiene y una si se detiene. Por tanto
\[
 U_e:\quad \#q_\infty^{-1}(v)=1
 \quad\Longleftrightarrow\quad T_e\downarrow.
\]
Un decisor uniforme para \(U_e\) decidiría \(\HALT\). Una teoría
recursivamente axiomatizable, correcta y completa para todos los \(U_e\)
permitiría decidirlos enumerando en paralelo pruebas de \(U_e\) y de su
negación. Éste es el núcleo del \emph{Teorema HMT de incompletitud uniforme
para multiplicidades proyectivas}; su demostración emplea una reducción
Gödel--Turing. La obstrucción uniforme no niega los selectores particulares
construidos sobre imágenes HMT calibradas.

\begin{controlbox}[title={Alcance del control de incompletitud uniforme}]
La reducción Gödel--Turing establece la obstrucción uniforme para las
multiplicidades proyectivas. No niega los selectores particulares construidos
sobre imágenes HMT calibradas ni convierte una obstrucción de decisión global
en una indeterminación de cada fibra concreta.
\end{controlbox}

\HMTRestoreHeadings

\section[Controles de la realización físico-dimensional]{Dominio y obstrucciones de las realizaciones física y dimensional}
\label{sec:int68-certificado-historico}

La realización físico-dimensional se somete a diez familias de control:
\begin{enumerate}
 \item concentración cilíndrica, masa uno y momentos finitos;
 \item espectros \((9,5)\) del bloque APP y \((7,3)\) del bloque de memoria;
 \item identidad \(\tanh^2(2\log\varphi)=5/9\);
 \item antiunitariedad en el modelo finito declarado;
 \item ley del producto semidirecto de Poincaré;
 \item identidades constitutivas del vacío;
 \item norma partida de los modos nulos y clausura material;
 \item selector \(\theta_*=(54/\pi)^2\) y punto fijo autodual;
 \item valores CGLMP de control;
 \item invariantes \((396,729,43,387)\) de la filtración nonádica y
       dominio \((81,56,13,324)\), con sector nulo separado, de la Ley
       General de Masas.
\end{enumerate}

\begin{controlbox}[title={Alcance exacto de los controles}]
Las diez familias comprueban identidades, dominios finitos y ejemplos
reproducibles. No sustituyen la funtorialidad completa de las cuatro
componentes, no construyen por sí solas el dominio común, no demuestran la
prolongación suave de Cartan ni promueven las interfaces físicas abiertas a
teoremas internos. La realización discreta conserva su demostración principal
en la \cref{sec:int68-realizacion-fisica-discreta}.
\end{controlbox}

\subsection{Dependencias functoriales de las realizaciones espectral, solenoidal, de calibre, cohomológica y determinantal}

Los controles anteriores sólo se aplican después de construir APP, TRIT, TPK,
el estado enriquecido común, la doble proyección y la recta dimensional. La
concentración cilíndrica y el espectro finito prueban propiedades de la
realización; no seleccionan retrospectivamente el estado. La periodicidad de
108 fases, el atlas (81/56/13/324), la involución constitutiva y los
programas cosmológicos conservan dominios y falsadores separados.

Las identidades de Dirac, von Neumann, CGLMP, Cartan y Poincaré se utilizan
como controles tipados del entrelazador correspondiente. Una coincidencia
numérica sin el mapa de dominio, el codominio y la acción sobre generadores no
constituye una realización física.

% Control científico del puente matricial. Las matrices de producción y la
% procedencia técnica permanecen en el expediente externo.
\chapter{Fidelidad genealógica del puente matricial: dominio, correspondencia y límites exactos}
\label{gmc:app:certificado}

\section{Identidades de naturalidad, positividad y compatibilidad del puente genealógico–matricial}

Los controles siguientes se obtienen reconstruyendo los operadores desde sus
datos de incidencia. Las identidades racionales se calculan exactamente; las
comparaciones de punto flotante declaran su cota de error.

\begin{longtable}{@{}L{.35\textwidth}L{.22\textwidth}L{.32\textwidth}@{}}
\toprule
Control & Estatuto & Error máximo o identidad\\
\midrule
\endfirsthead
\toprule
Control & Estatuto & Error máximo o identidad\\
\midrule
\endhead
cuatro \PVM qutrit & verificado & \(4.72\times10^{-16}\)\\
idempotencia de los doce proyectores & verificada & \(3.34\times10^{-16}\)\\
solapamientos MUB & verificados & \(\lvert\Tr(PQ)-1/3\rvert<3.34\times10^{-16}\)\\
obstrucción de columna QLS & exacta & mínimo \(1/3\)\\
QMS aditivo \((9,3)\) & verificado & filas y columnas exactas a tolerancia\\
QMS de tres contextos \((9,3)\) & verificado & error \(<4.21\times10^{-16}\)\\
no conmutatividad de celdas & verificada & norma máxima \(0.074074\ldots\)\\
QMS atlas \((12,3)\) & verificado & error \(<4.24\times10^{-16}\)\\
descomposición semiclásica & exacta & error de reconstrucción cero\\
cubo qutrit de orden nueve & verificado & error de línea \(<4.21\times10^{-16}\)\\
cubo qutrit de orden doce & verificado & error de línea \(<4.24\times10^{-16}\)\\
\(\mathbb P^1(\mathbb Z/9)\to\mathbb P^1(\mathbb F_3)\) & exacto &
\(12\to4\) con fibras \((3,3,3,3)\)\\
truncación--compresión & exacta & conmutador cero\\
lector directo \(1/\pi\) & exacto & normalización unitaria\\
calibre MUB \(3/\pi\) & verificado & probabilidad \(1/\pi\), error \(<3.85\times10^{-16}\)\\
polinomio de Gram APP & exacto & \(t^5(t-1)(t-5/7)(t-1/3)^2\)\\
cuatro intersecciones de Halmos & exactas & \((64,1,5,5)\)\\
álgebra \(C^*\) de APP en \(d=2\) & exacta &
\(\mathbb C^4\oplus M_2\oplus M_2\), dimensión \(12\)\\
\bottomrule
\end{longtable}

La tolerancia de los controles numéricos es \(2\times10^{-10}\). Las
afirmaciones racionales no usan tolerancia: se obtienen por fracciones exactas
y eliminación gaussiana sobre \(\mathbb Q\).

\Needspace{30\baselineskip}
\section{Coeficientes exactos del polinomio}

La construcción produce
\[
 \det(tI-CC^*)
 =t^9-\frac{50}{21}t^8+\frac{124}{63}t^7
  -\frac23t^6+\frac5{63}t^5,
 \tag{GMC-A.1}
\]
y coincide coeficiente a coeficiente con
\[
 t^5(t-1)(t-5/7)(t-1/3)^2.
 \tag{GMC-A.2}
\]
La comprobación construye \(CC^*\) desde los cardinales exactos de las
intersecciones de fibras y aplica Faddeev--LeVerrier sobre \(\mathbb Q\); no
depende de una diagonalización numérica.

\section{Obstrucciones a la identificación de equivalencias no demostradas en el puente matricial}

El certificado finito no sustituye la filtración nonádica a profundidad arbitraria, la
Ley General de Masas ni las derivaciones de Planck. Tampoco demuestra por sí
solo la fidelidad APP--TRIT--TPK de todo QMS cíclico, la equivarianza completa
de la calibración \(\Xi\), la identificación de la memoria TPK con un sistema
auxiliar de Stinespring, la divisibilidad infinita de un canal MPS ni una
universalidad categórica. Cada una de esas afirmaciones conserva su residencia
principal y sus propios falsadores.



\chapter{Dominio de semillas, emisiones y expansiones verificables}
\label{app:semillas-emisiones-expansiones-acumulativas}

Este bloque conserva los dos niveles que sostienen la publicación numérica:
el censo exhaustivo de condiciones iniciales y emisiones del TPK, y las
expansiones extensas producidas por los lectores racionales. El catálogo
completo fija el paso
\[
104\,976\longrightarrow468\longrightarrow243,
\]
mientras que las expansiones de diez mil cifras documentan la prolongación
efectiva de las secciones de \(\pi\), \(e\), \(\varphi\) y
\(\alpha_{\rm HMT}\).

El catálogo de las (468) emisiones y las (243) palabras visibles se publica
íntegramente a continuación; las cuatro expansiones de diez mil cifras se
publican íntegramente en la genealogía de constantes. Esta sección terminal
enlaza ambos desarrollos sin reducir sus demostraciones, tablas ni etiquetas.

\begin{HMTScientificOwnerSubordinated}
\chapter{Dominio de semillas APP y catálogo canónico de emisiones}
\label{app:catalogo-semillas-app-rev9}
\index{semilla APP}
\index{catálogo de emisiones TPK}

\section{Dominio exhaustivo de condiciones iniciales}

Sea
\[
 I_9=\{0,1,\ldots,8\},
 \qquad D=\{N,E,S,O\},
 \qquad \mathcal S=I_9^2\times D.
\]
Un elemento de \(\mathcal S\) fija la posición y la orientación iniciales de
un cursor. El emisor mínimo utiliza dos cursores sobre el mismo soporte APP:
uno para la evaluación aditiva y otro para la evaluación multiplicativa. Su
dominio completo es
\[
 \Omega_{\rm seed}
 =\mathcal S_{+}\times\mathcal S_{\times}
 \cong\mathcal S^2,
 \qquad
 |\mathcal S|=9^2\cdot4=324,
 \qquad
 |\Omega_{\rm seed}|=324^2=104\,976.
\]

\begin{definition}[Semilla completa del emisor APP--TPK]
Una semilla es una pareja ordenada
\[
 \omega=(s_+,s_\times),
 \qquad
 s_+=(i_+,j_+,d_+),\quad
 s_\times=(i_\times,j_\times,d_\times),
 \qquad s_+,s_\times\in\mathcal S.
\]
Las \(104\,976\) semillas son exactamente los elementos de
\(\Omega_{\rm seed}\). Las filas del catálogo pertenecen al tipo emisión y
cada una representa una fibra de este dominio de semillas.
\end{definition}

La recurrencia de cincuenta y cuatro pasos define la aplicación finita
\[
 T_{\min}:\Omega_{\rm seed}\longrightarrow\mathcal U_6^{\rm cat},
 \qquad
 \omega\longmapsto U_6=(U_1,\ldots,U_6).
\]
Su imagen contiene \(468\) sextetos decimales distintos. La reducción
componente a componente produce después \(243\) palabras distintas en el
ambiente \(\mathbb F_3^6\), de cardinal \(729\). Por tanto, los tres niveles
quedan tipados como
\[
 \underbrace{104\,976}_{\text{semillas }\omega}
 \xrightarrow{\;T_{\min}\;}
 \underbrace{468}_{\text{emisiones }U_6}
 \xrightarrow{\;\rho_3\;}
 \underbrace{243}_{\text{palabras visibles}}
 \subset
 \underbrace{\mathbb F_3^6}_{729\text{ palabras}}.
\]

\begin{theorem}[Censo exhaustivo del dominio y de la imagen]
El catálogo canónico contiene \(468\) emisiones \(U_6\) distintas y
\(\#T_{\min}^{-1}(U_6)\) es la cardinalidad de la fibra
\(T_{\min}^{-1}(U_6)\). Sus multiplicidades satisfacen
\[
 432\cdot144+18\cdot432+18\cdot1944=104\,976,
\]
de modo que la suma de todas las fibras recupera exactamente el dominio
\(\Omega_{\rm seed}\). La proyección directa
\(\rho_3(U_6)\) tiene \(243\) valores distintos.
\end{theorem}

\paragraph{Demostración del censo exhaustivo \(104\,976\to468\to243\) y de las multiplicidades de las fibras de \(T_{\min}\)}
La recurrencia que define \(T_{\min}\) reside en
\cref{pf84:E02-16}; la factorización de su imagen y el cómputo de las fibras,
en \cref{pf84:E02-20,cor:censo-proyeccion-observable}. Este apéndice no
introduce una segunda enumeración: materializa, con todos sus campos, el
catálogo canónico producido por esos resultados.

\section{Catálogo completo de las 468 emisiones}

La tabla siguiente incorpora los diez campos del catálogo exhaustivo. El
ordinal de la primera columna fija un orden de lectura, con un tipo distinto
del de las semillas y de los identificadores matemáticos. Las columnas
abreviadas conservan
\texttt{diag\_c}, \texttt{hplus\_type}, \texttt{px\_size},
\texttt{px\_res\_type}, \texttt{delta\_diag\_subset},
    \texttt{delta\_orbit}, \texttt{Esig} y \(\rho_3(U_6)\).

\begingroup
\tiny
\setlength{\tabcolsep}{2pt}
\renewcommand{\arraystretch}{1.12}
\begin{longtable}{@{}r
  >{\ttfamily\raggedright\arraybackslash}p{.245\textwidth}
  >{\raggedright\arraybackslash}p{.085\textwidth}
  >{\ttfamily\raggedright\arraybackslash}p{.105\textwidth}
  >{\ttfamily\raggedright\arraybackslash}p{.20\textwidth}
  >{\ttfamily\raggedright\arraybackslash}p{.15\textwidth}@{}}
\caption{Catálogo canónico completo de emisiones U6.}
\label{tab:catalogo-emisiones-u6-rev9}\\
\toprule
\# & \normalfont\(U_6\) / \(\rho_3(U_6)\)
   & \normalfont\(\#T_{\min}^{-1}(U_6)\) / \texttt{diag\_c}
   & \normalfont\texttt{hplus\_type} / \texttt{px\_size} / \texttt{px\_res\_type}
   & \normalfont\texttt{delta\_diag\_subset} / \texttt{delta\_orbit}
   & \normalfont\texttt{Esig}\\
\midrule
\endfirsthead
\multicolumn{6}{c}{\tablename\ \thetable\ (continuación)}\\
\toprule
\# & \normalfont\(U_6\) / \(\rho_3(U_6)\)
   & \normalfont\(\#T_{\min}^{-1}(U_6)\) / \texttt{diag\_c}
   & \normalfont\texttt{hplus\_type} / \texttt{px\_size} / \texttt{px\_res\_type}
   & \normalfont\texttt{delta\_diag\_subset} / \texttt{delta\_orbit}
   & \normalfont\texttt{Esig}\\
\midrule
\endhead
\midrule
\multicolumn{6}{r}{Continúa en la página siguiente.}\\
\endfoot
\bottomrule
\endlastfoot
1 & 252|109|252|903|850|850 / 010011 & 144 / 0 & NO / 6 / A & 0,1,4,5 / 0,1,4,5 & 3,9,3,1,7,7 \\
2 & 232|189|292|923|870|810 / 101200 & 144 / 0 & NO / 6 / A & 0,1,4,5 / 0,1,4,5 & 1,7,7,3,9,3 \\
3 & 923|870|810|232|189|292 / 200101 & 144 / 0 & ES / 6 / A & 0,1,4,5 / 0,1,4,5 & 3,9,3,1,7,7 \\
4 & 903|850|850|252|109|252 / 011010 & 144 / 0 & ES / 6 / A & 0,1,4,5 / 0,1,4,5 & 1,7,7,3,9,3 \\
5 & 272|169|272|923|810|870 / 212200 & 144 / 0 & NO / 8 / A & 0,1,4,5 / 0,1,4,5 & 5,5,5,3,3,9 \\
6 & 252|149|212|943|830|830 / 022122 & 144 / 0 & NO / 8 / A & 0,1,4,5 / 0,1,4,5 & 3,3,9,5,5,5 \\
7 & 943|830|830|252|149|212 / 122022 & 144 / 0 & ES / 8 / A & 0,1,4,5 / 0,1,4,5 & 5,5,5,3,3,9 \\
8 & 923|810|870|272|169|272 / 200212 & 144 / 0 & ES / 8 / A & 0,1,4,5 / 0,1,4,5 & 3,3,9,5,5,5 \\
9 & 227|114|227|998|885|885 / 202200 & 1944 / 0 & NO / 45 / H & 0,1,2,3,4,5,6,7,8 / 0,1,2,3,4,5,6,7,8 & 0,0,0,0,0,0 \\
10 & 292|189|232|903|850|850 / 101011 & 144 / 0 & NO / 8 / B & 2,3,6,8 / 0,1,4,6 & 7,7,1,1,7,7 \\
11 & 232|189|292|963|850|890 / 101012 & 144 / 0 & NO / 8 / B & 2,3,6,8 / 0,1,4,6 & 1,7,7,7,7,1 \\
12 & 998|885|885|227|114|227 / 200202 & 1944 / 0 & ES / 45 / H & 0,1,2,3,4,5,6,7,8 / 0,1,2,3,4,5,6,7,8 & 0,0,0,0,0,0 \\
13 & 963|850|890|232|189|292 / 012101 & 144 / 0 & ES / 8 / B & 2,3,6,8 / 0,1,4,6 & 7,7,1,1,7,7 \\
14 & 903|850|850|292|189|232 / 011101 & 144 / 0 & ES / 8 / B & 2,3,6,8 / 0,1,4,6 & 1,7,7,7,7,1 \\
15 & 292|129|292|943|830|830 / 101122 & 144 / 0 & NO / 8 / A & 1,2,3,4 / 0,1,2,3 & 7,1,7,5,5,5 \\
16 & 292|189|232|923|810|870 / 101200 & 144 / 0 & NO / 6 / A & 2,3,7 / 0,1,5 & 7,7,1,3,3,9 \\
17 & 272|169|272|963|890|850 / 212021 & 144 / 0 & NO / 8 / A & 1,2,3,4 / 0,1,2,3 & 5,5,5,7,1,7 \\
18 & 252|149|212|963|850|890 / 022012 & 144 / 0 & NO / 6 / A & 2,3,7 / 0,1,5 & 3,3,9,7,7,1 \\
19 & 963|890|850|272|169|272 / 021212 & 144 / 0 & ES / 8 / A & 1,2,3,4 / 0,1,2,3 & 7,1,7,5,5,5 \\
20 & 963|850|890|252|149|212 / 012022 & 144 / 0 & ES / 6 / A & 2,3,7 / 0,1,5 & 7,7,1,3,3,9 \\
21 & 943|830|830|292|129|292 / 122101 & 144 / 0 & ES / 8 / A & 1,2,3,4 / 0,1,2,3 & 5,5,5,7,1,7 \\
22 & 923|810|870|292|189|232 / 200101 & 144 / 0 & ES / 6 / A & 2,3,7 / 0,1,5 & 3,3,9,7,7,1 \\
23 & 292|189|232|943|830|830 / 101122 & 144 / 0 & NO / 6 / A & 1,4,6,8 / 0,2,4,6 & 7,7,1,5,5,5 \\
24 & 212|149|252|963|890|850 / 220021 & 144 / 0 & NO / 8 / A & 1,4,6,8 / 0,2,4,6 & 9,3,3,7,1,7 \\
25 & 272|169|272|963|850|890 / 212012 & 144 / 0 & NO / 6 / A & 1,4,6,8 / 0,2,4,6 & 5,5,5,7,7,1 \\
26 & 292|129|292|983|810|810 / 101200 & 144 / 0 & NO / 8 / A & 1,4,6,8 / 0,2,4,6 & 7,1,7,9,3,3 \\
27 & 963|850|890|272|169|272 / 012212 & 144 / 0 & ES / 6 / A & 1,4,6,8 / 0,2,4,6 & 7,7,1,5,5,5 \\
28 & 983|810|810|292|129|292 / 200101 & 144 / 0 & ES / 8 / A & 1,4,6,8 / 0,2,4,6 & 9,3,3,7,1,7 \\
29 & 943|830|830|292|189|232 / 122101 & 144 / 0 & ES / 6 / A & 1,4,6,8 / 0,2,4,6 & 5,5,5,7,7,1 \\
30 & 963|890|850|212|149|252 / 021220 & 144 / 0 & ES / 8 / A & 1,4,6,8 / 0,2,4,6 & 7,1,7,9,3,3 \\
31 & 292|129|292|963|890|850 / 101021 & 144 / 0 & NO / 4 / B & 0,5 / 0,4 & 7,1,7,7,1,7 \\
32 & 963|890|850|292|129|292 / 021101 & 144 / 0 & ES / 4 / B & 0,5 / 0,4 & 7,1,7,7,1,7 \\
33 & 272|169|272|983|810|810 / 212200 & 144 / 0 & NO / 6 / A & 6,7,8 / 0,1,2 & 5,5,5,9,3,3 \\
34 & 212|149|252|943|830|830 / 220122 & 144 / 0 & NO / 6 / A & 6,7,8 / 0,1,2 & 9,3,3,5,5,5 \\
35 & 943|830|830|212|149|252 / 122220 & 144 / 0 & ES / 6 / A & 6,7,8 / 0,1,2 & 5,5,5,9,3,3 \\
36 & 983|810|810|272|169|272 / 200212 & 144 / 0 & ES / 6 / A & 6,7,8 / 0,1,2 & 9,3,3,5,5,5 \\
37 & 252|109|252|943|830|830 / 010122 & 144 / 0 & NO / 6 / A & 1,2,3,4 / 0,1,2,3 & 3,9,3,5,5,5 \\
38 & 272|169|272|923|870|810 / 212200 & 144 / 0 & NO / 6 / A & 1,2,3,4 / 0,1,2,3 & 5,5,5,3,9,3 \\
39 & 923|870|810|272|169|272 / 200212 & 144 / 0 & ES / 6 / A & 1,2,3,4 / 0,1,2,3 & 3,9,3,5,5,5 \\
40 & 943|830|830|252|109|252 / 122010 & 144 / 0 & ES / 6 / A & 1,2,3,4 / 0,1,2,3 & 5,5,5,3,9,3 \\
41 & 252|109|252|923|870|810 / 010200 & 144 / 0 & NO / 4 / B & 2,3 / 0,1 & 3,9,3,3,9,3 \\
42 & 923|870|810|252|109|252 / 200010 & 144 / 0 & ES / 4 / B & 2,3 / 0,1 & 3,9,3,3,9,3 \\
43 & 272|169|272|903|850|850 / 212011 & 144 / 0 & NO / 6 / A & 0,5,7 / 0,2,4 & 5,5,5,1,7,7 \\
44 & 232|189|292|943|830|830 / 101122 & 144 / 0 & NO / 6 / A & 0,5,7 / 0,2,4 & 1,7,7,5,5,5 \\
45 & 943|830|830|232|189|292 / 122101 & 144 / 0 & ES / 6 / A & 0,5,7 / 0,2,4 & 5,5,5,1,7,7 \\
46 & 903|850|850|272|169|272 / 011212 & 144 / 0 & ES / 6 / A & 0,5,7 / 0,2,4 & 1,7,7,5,5,5 \\
47 & 212|149|252|923|810|870 / 220200 & 144 / 0 & NO / 8 / B & 0,5,6,8 / 0,1,3,4 & 9,3,3,3,3,9 \\
48 & 252|149|212|983|810|810 / 022200 & 144 / 0 & NO / 8 / B & 0,5,6,8 / 0,1,3,4 & 3,3,9,9,3,3 \\
49 & 983|810|810|252|149|212 / 200022 & 144 / 0 & ES / 8 / B & 0,5,6,8 / 0,1,3,4 & 9,3,3,3,3,9 \\
50 & 923|810|870|212|149|252 / 200220 & 144 / 0 & ES / 8 / B & 0,5,6,8 / 0,1,3,4 & 3,3,9,9,3,3 \\
51 & 272|169|272|943|830|830 / 212122 & 432 / 0 & NO / 12 / B & 0,2,3,5,6,8 / 0,1,3,4,6,7 & 5,5,5,5,5,5 \\
52 & 943|830|830|272|169|272 / 122212 & 432 / 0 & ES / 12 / B & 0,2,3,5,6,8 / 0,1,3,4,6,7 & 5,5,5,5,5,5 \\
53 & 694|438|581|232|189|292 / 102101 & 144 / 1 & NO / 6 / A & 0,3,4,8 / 0,1,4,5 & 3,9,3,1,7,7 \\
54 & 674|418|521|252|109|252 / 212010 & 144 / 1 & NO / 6 / A & 0,3,4,8 / 0,1,4,5 & 1,7,7,3,9,3 \\
55 & 252|109|252|674|418|521 / 010212 & 144 / 1 & ES / 6 / A & 0,3,4,8 / 0,1,4,5 & 3,9,3,1,7,7 \\
56 & 232|189|292|694|438|581 / 101102 & 144 / 1 & ES / 6 / A & 0,3,4,8 / 0,1,4,5 & 1,7,7,3,9,3 \\
57 & 614|498|501|252|149|212 / 200022 & 144 / 1 & NO / 8 / A & 0,3,4,8 / 0,1,4,5 & 5,5,5,3,3,9 \\
58 & 694|478|541|272|169|272 / 111212 & 144 / 1 & NO / 8 / A & 0,3,4,8 / 0,1,4,5 & 3,3,9,5,5,5 \\
59 & 272|169|272|694|478|541 / 212111 & 144 / 1 & ES / 8 / A & 0,3,4,8 / 0,1,4,5 & 5,5,5,3,3,9 \\
60 & 252|149|212|614|498|501 / 022200 & 144 / 1 & ES / 8 / A & 0,3,4,8 / 0,1,4,5 & 3,3,9,5,5,5 \\
61 & 669|443|556|227|114|227 / 021202 & 1944 / 1 & NO / 45 / H & 0,1,2,3,4,5,6,7,8 / 0,1,2,3,4,5,6,7,8 & 0,0,0,0,0,0 \\
62 & 634|418|561|232|189|292 / 110101 & 144 / 1 & NO / 8 / B & 1,2,5,7 / 0,1,4,6 & 7,7,1,1,7,7 \\
63 & 674|418|521|292|189|232 / 212101 & 144 / 1 & NO / 8 / B & 1,2,5,7 / 0,1,4,6 & 1,7,7,7,7,1 \\
64 & 227|114|227|669|443|556 / 202021 & 1944 / 1 & ES / 45 / H & 0,1,2,3,4,5,6,7,8 / 0,1,2,3,4,5,6,7,8 & 0,0,0,0,0,0 \\
65 & 292|189|232|674|418|521 / 101212 & 144 / 1 & ES / 8 / B & 1,2,5,7 / 0,1,4,6 & 7,7,1,1,7,7 \\
66 & 232|189|292|634|418|561 / 101110 & 144 / 1 & ES / 8 / B & 1,2,5,7 / 0,1,4,6 & 1,7,7,7,7,1 \\
67 & 634|458|521|272|169|272 / 122212 & 144 / 1 & NO / 8 / A & 0,1,2,3 / 0,1,2,3 & 7,1,7,5,5,5 \\
68 & 634|418|561|252|149|212 / 110022 & 144 / 1 & NO / 6 / A & 1,2,6 / 0,1,5 & 7,7,1,3,3,9 \\
69 & 614|498|501|292|129|292 / 200101 & 144 / 1 & NO / 8 / A & 0,1,2,3 / 0,1,2,3 & 5,5,5,7,1,7 \\
70 & 694|478|541|292|189|232 / 111101 & 144 / 1 & NO / 6 / A & 1,2,6 / 0,1,5 & 3,3,9,7,7,1 \\
71 & 292|129|292|614|498|501 / 101200 & 144 / 1 & ES / 8 / A & 0,1,2,3 / 0,1,2,3 & 7,1,7,5,5,5 \\
72 & 292|189|232|694|478|541 / 101111 & 144 / 1 & ES / 6 / A & 1,2,6 / 0,1,5 & 7,7,1,3,3,9 \\
73 & 272|169|272|634|458|521 / 212122 & 144 / 1 & ES / 8 / A & 0,1,2,3 / 0,1,2,3 & 5,5,5,7,1,7 \\
74 & 252|149|212|634|418|561 / 022110 & 144 / 1 & ES / 6 / A & 1,2,6 / 0,1,5 & 3,3,9,7,7,1 \\
75 & 634|418|561|272|169|272 / 110212 & 144 / 1 & NO / 6 / A & 0,3,5,7 / 0,2,4,6 & 7,7,1,5,5,5 \\
76 & 654|478|581|292|129|292 / 012101 & 144 / 1 & NO / 8 / A & 0,3,5,7 / 0,2,4,6 & 9,3,3,7,1,7 \\
77 & 614|498|501|292|189|232 / 200101 & 144 / 1 & NO / 6 / A & 0,3,5,7 / 0,2,4,6 & 5,5,5,7,7,1 \\
78 & 634|458|521|212|149|252 / 122220 & 144 / 1 & NO / 8 / A & 0,3,5,7 / 0,2,4,6 & 7,1,7,9,3,3 \\
79 & 292|189|232|614|498|501 / 101200 & 144 / 1 & ES / 6 / A & 0,3,5,7 / 0,2,4,6 & 7,7,1,5,5,5 \\
80 & 212|149|252|634|458|521 / 220122 & 144 / 1 & ES / 8 / A & 0,3,5,7 / 0,2,4,6 & 9,3,3,7,1,7 \\
81 & 272|169|272|634|418|561 / 212110 & 144 / 1 & ES / 6 / A & 0,3,5,7 / 0,2,4,6 & 5,5,5,7,7,1 \\
82 & 292|129|292|654|478|581 / 101012 & 144 / 1 & ES / 8 / A & 0,3,5,7 / 0,2,4,6 & 7,1,7,9,3,3 \\
83 & 634|458|521|292|129|292 / 122101 & 144 / 1 & NO / 4 / B & 4,8 / 0,4 & 7,1,7,7,1,7 \\
84 & 292|129|292|634|458|521 / 101122 & 144 / 1 & ES / 4 / B & 4,8 / 0,4 & 7,1,7,7,1,7 \\
85 & 614|498|501|212|149|252 / 200220 & 144 / 1 & NO / 6 / A & 5,6,7 / 0,1,2 & 5,5,5,9,3,3 \\
86 & 654|478|581|272|169|272 / 012212 & 144 / 1 & NO / 6 / A & 5,6,7 / 0,1,2 & 9,3,3,5,5,5 \\
87 & 272|169|272|654|478|581 / 212012 & 144 / 1 & ES / 6 / A & 5,6,7 / 0,1,2 & 5,5,5,9,3,3 \\
88 & 212|149|252|614|498|501 / 220200 & 144 / 1 & ES / 6 / A & 5,6,7 / 0,1,2 & 9,3,3,5,5,5 \\
89 & 694|438|581|272|169|272 / 102212 & 144 / 1 & NO / 6 / A & 0,1,2,3 / 0,1,2,3 & 3,9,3,5,5,5 \\
90 & 614|498|501|252|109|252 / 200010 & 144 / 1 & NO / 6 / A & 0,1,2,3 / 0,1,2,3 & 5,5,5,3,9,3 \\
91 & 252|109|252|614|498|501 / 010200 & 144 / 1 & ES / 6 / A & 0,1,2,3 / 0,1,2,3 & 3,9,3,5,5,5 \\
92 & 272|169|272|694|438|581 / 212102 & 144 / 1 & ES / 6 / A & 0,1,2,3 / 0,1,2,3 & 5,5,5,3,9,3 \\
93 & 694|438|581|252|109|252 / 102010 & 144 / 1 & NO / 4 / B & 1,2 / 0,1 & 3,9,3,3,9,3 \\
94 & 252|109|252|694|438|581 / 010102 & 144 / 1 & ES / 4 / B & 1,2 / 0,1 & 3,9,3,3,9,3 \\
95 & 614|498|501|232|189|292 / 200101 & 144 / 1 & NO / 6 / A & 4,6,8 / 0,2,4 & 5,5,5,1,7,7 \\
96 & 674|418|521|272|169|272 / 212212 & 144 / 1 & NO / 6 / A & 4,6,8 / 0,2,4 & 1,7,7,5,5,5 \\
97 & 272|169|272|674|418|521 / 212212 & 144 / 1 & ES / 6 / A & 4,6,8 / 0,2,4 & 5,5,5,1,7,7 \\
98 & 232|189|292|614|498|501 / 101200 & 144 / 1 & ES / 6 / A & 4,6,8 / 0,2,4 & 1,7,7,5,5,5 \\
99 & 654|478|581|252|149|212 / 012022 & 144 / 1 & NO / 8 / B & 4,5,7,8 / 0,1,3,4 & 9,3,3,3,3,9 \\
100 & 694|478|541|212|149|252 / 111220 & 144 / 1 & NO / 8 / B & 4,5,7,8 / 0,1,3,4 & 3,3,9,9,3,3 \\
101 & 212|149|252|694|478|541 / 220111 & 144 / 1 & ES / 8 / B & 4,5,7,8 / 0,1,3,4 & 9,3,3,3,3,9 \\
102 & 252|149|212|654|478|581 / 022012 & 144 / 1 & ES / 8 / B & 4,5,7,8 / 0,1,3,4 & 3,3,9,9,3,3 \\
103 & 614|498|501|272|169|272 / 200212 & 432 / 1 & NO / 12 / B & 1,2,4,5,7,8 / 0,1,3,4,6,7 & 5,5,5,5,5,5 \\
104 & 272|169|272|614|498|501 / 212200 & 432 / 1 & ES / 12 / B & 1,2,4,5,7,8 / 0,1,3,4,6,7 & 5,5,5,5,5,5 \\
105 & 923|870|810|561|418|634 / 200011 & 144 / 2 & NO / 6 / A & 2,3,7,8 / 0,1,4,5 & 3,9,3,1,7,7 \\
106 & 903|850|850|581|438|694 / 011201 & 144 / 2 & NO / 6 / A & 2,3,7,8 / 0,1,4,5 & 1,7,7,3,9,3 \\
107 & 581|438|694|903|850|850 / 201011 & 144 / 2 & ES / 6 / A & 2,3,7,8 / 0,1,4,5 & 3,9,3,1,7,7 \\
108 & 561|418|634|923|870|810 / 011200 & 144 / 2 & ES / 6 / A & 2,3,7,8 / 0,1,4,5 & 1,7,7,3,9,3 \\
109 & 943|830|830|581|478|654 / 122210 & 144 / 2 & NO / 8 / A & 2,3,7,8 / 0,1,4,5 & 5,5,5,3,3,9 \\
110 & 923|810|870|501|498|614 / 200002 & 144 / 2 & NO / 8 / A & 2,3,7,8 / 0,1,4,5 & 3,3,9,5,5,5 \\
111 & 501|498|614|923|810|870 / 002200 & 144 / 2 & ES / 8 / A & 2,3,7,8 / 0,1,4,5 & 5,5,5,3,3,9 \\
112 & 581|478|654|943|830|830 / 210122 & 144 / 2 & ES / 8 / A & 2,3,7,8 / 0,1,4,5 & 3,3,9,5,5,5 \\
113 & 998|885|885|556|443|669 / 200120 & 1944 / 2 & NO / 45 / H & 0,1,2,3,4,5,6,7,8 / 0,1,2,3,4,5,6,7,8 & 0,0,0,0,0,0 \\
114 & 963|850|890|561|418|634 / 012011 & 144 / 2 & NO / 8 / B & 0,1,4,6 / 0,1,4,6 & 7,7,1,1,7,7 \\
115 & 903|850|850|521|418|674 / 011212 & 144 / 2 & NO / 8 / B & 0,1,4,6 / 0,1,4,6 & 1,7,7,7,7,1 \\
116 & 556|443|669|998|885|885 / 120200 & 1944 / 2 & ES / 45 / H & 0,1,2,3,4,5,6,7,8 / 0,1,2,3,4,5,6,7,8 & 0,0,0,0,0,0 \\
117 & 521|418|674|903|850|850 / 212011 & 144 / 2 & ES / 8 / B & 0,1,4,6 / 0,1,4,6 & 7,7,1,1,7,7 \\
118 & 561|418|634|963|850|890 / 011012 & 144 / 2 & ES / 8 / B & 0,1,4,6 / 0,1,4,6 & 1,7,7,7,7,1 \\
119 & 963|890|850|501|498|614 / 021002 & 144 / 2 & NO / 8 / A & 0,1,2,8 / 0,1,2,3 & 7,1,7,5,5,5 \\
120 & 963|850|890|581|478|654 / 012210 & 144 / 2 & NO / 6 / A & 0,1,5 / 0,1,5 & 7,7,1,3,3,9 \\
121 & 943|830|830|521|458|634 / 122221 & 144 / 2 & NO / 8 / A & 0,1,2,8 / 0,1,2,3 & 5,5,5,7,1,7 \\
122 & 923|810|870|521|418|674 / 200212 & 144 / 2 & NO / 6 / A & 0,1,5 / 0,1,5 & 3,3,9,7,7,1 \\
123 & 521|458|634|943|830|830 / 221122 & 144 / 2 & ES / 8 / A & 0,1,2,8 / 0,1,2,3 & 7,1,7,5,5,5 \\
124 & 521|418|674|923|810|870 / 212200 & 144 / 2 & ES / 6 / A & 0,1,5 / 0,1,5 & 7,7,1,3,3,9 \\
125 & 501|498|614|963|890|850 / 002021 & 144 / 2 & ES / 8 / A & 0,1,2,8 / 0,1,2,3 & 5,5,5,7,1,7 \\
126 & 581|478|654|963|850|890 / 210012 & 144 / 2 & ES / 6 / A & 0,1,5 / 0,1,5 & 3,3,9,7,7,1 \\
127 & 963|850|890|501|498|614 / 012002 & 144 / 2 & NO / 6 / A & 2,4,6,8 / 0,2,4,6 & 7,7,1,5,5,5 \\
128 & 983|810|810|521|458|634 / 200221 & 144 / 2 & NO / 8 / A & 2,4,6,8 / 0,2,4,6 & 9,3,3,7,1,7 \\
129 & 943|830|830|521|418|674 / 122212 & 144 / 2 & NO / 6 / A & 2,4,6,8 / 0,2,4,6 & 5,5,5,7,7,1 \\
130 & 963|890|850|541|478|694 / 021111 & 144 / 2 & NO / 8 / A & 2,4,6,8 / 0,2,4,6 & 7,1,7,9,3,3 \\
131 & 521|418|674|943|830|830 / 212122 & 144 / 2 & ES / 6 / A & 2,4,6,8 / 0,2,4,6 & 7,7,1,5,5,5 \\
132 & 541|478|694|963|890|850 / 111021 & 144 / 2 & ES / 8 / A & 2,4,6,8 / 0,2,4,6 & 9,3,3,7,1,7 \\
133 & 501|498|614|963|850|890 / 002012 & 144 / 2 & ES / 6 / A & 2,4,6,8 / 0,2,4,6 & 5,5,5,7,7,1 \\
134 & 521|458|634|983|810|810 / 221200 & 144 / 2 & ES / 8 / A & 2,4,6,8 / 0,2,4,6 & 7,1,7,9,3,3 \\
135 & 963|890|850|521|458|634 / 021221 & 144 / 2 & NO / 4 / B & 3,7 / 0,4 & 7,1,7,7,1,7 \\
136 & 521|458|634|963|890|850 / 221021 & 144 / 2 & ES / 4 / B & 3,7 / 0,4 & 7,1,7,7,1,7 \\
137 & 943|830|830|541|478|694 / 122111 & 144 / 2 & NO / 6 / A & 4,5,6 / 0,1,2 & 5,5,5,9,3,3 \\
138 & 983|810|810|501|498|614 / 200002 & 144 / 2 & NO / 6 / A & 4,5,6 / 0,1,2 & 9,3,3,5,5,5 \\
139 & 501|498|614|983|810|810 / 002200 & 144 / 2 & ES / 6 / A & 4,5,6 / 0,1,2 & 5,5,5,9,3,3 \\
140 & 541|478|694|943|830|830 / 111122 & 144 / 2 & ES / 6 / A & 4,5,6 / 0,1,2 & 9,3,3,5,5,5 \\
141 & 923|870|810|501|498|614 / 200002 & 144 / 2 & NO / 6 / A & 0,1,2,8 / 0,1,2,3 & 3,9,3,5,5,5 \\
142 & 943|830|830|581|438|694 / 122201 & 144 / 2 & NO / 6 / A & 0,1,2,8 / 0,1,2,3 & 5,5,5,3,9,3 \\
143 & 581|438|694|943|830|830 / 201122 & 144 / 2 & ES / 6 / A & 0,1,2,8 / 0,1,2,3 & 3,9,3,5,5,5 \\
144 & 501|498|614|923|870|810 / 002200 & 144 / 2 & ES / 6 / A & 0,1,2,8 / 0,1,2,3 & 5,5,5,3,9,3 \\
145 & 923|870|810|581|438|694 / 200201 & 144 / 2 & NO / 4 / B & 0,1 / 0,1 & 3,9,3,3,9,3 \\
146 & 581|438|694|923|870|810 / 201200 & 144 / 2 & ES / 4 / B & 0,1 / 0,1 & 3,9,3,3,9,3 \\
147 & 943|830|830|561|418|634 / 122011 & 144 / 2 & NO / 6 / A & 3,5,7 / 0,2,4 & 5,5,5,1,7,7 \\
148 & 903|850|850|501|498|614 / 011002 & 144 / 2 & NO / 6 / A & 3,5,7 / 0,2,4 & 1,7,7,5,5,5 \\
149 & 501|498|614|903|850|850 / 002011 & 144 / 2 & ES / 6 / A & 3,5,7 / 0,2,4 & 5,5,5,1,7,7 \\
150 & 561|418|634|943|830|830 / 011122 & 144 / 2 & ES / 6 / A & 3,5,7 / 0,2,4 & 1,7,7,5,5,5 \\
151 & 983|810|810|581|478|654 / 200210 & 144 / 2 & NO / 8 / B & 3,4,6,7 / 0,1,3,4 & 9,3,3,3,3,9 \\
152 & 923|810|870|541|478|694 / 200111 & 144 / 2 & NO / 8 / B & 3,4,6,7 / 0,1,3,4 & 3,3,9,9,3,3 \\
153 & 541|478|694|923|810|870 / 111200 & 144 / 2 & ES / 8 / B & 3,4,6,7 / 0,1,3,4 & 9,3,3,3,3,9 \\
154 & 581|478|654|983|810|810 / 210200 & 144 / 2 & ES / 8 / B & 3,4,6,7 / 0,1,3,4 & 3,3,9,9,3,3 \\
155 & 943|830|830|501|498|614 / 122002 & 432 / 2 & NO / 12 / B & 0,1,3,4,6,7 / 0,1,3,4,6,7 & 5,5,5,5,5,5 \\
156 & 501|498|614|943|830|830 / 002122 & 432 / 2 & ES / 12 / B & 0,1,3,4,6,7 / 0,1,3,4,6,7 & 5,5,5,5,5,5 \\
157 & 252|212|149|890|850|963 / 022210 & 144 / 3 & NO / 6 / A & 1,2,6,7 / 0,1,4,5 & 3,9,3,1,7,7 \\
158 & 232|292|189|810|870|923 / 110002 & 144 / 3 & NO / 6 / A & 1,2,6,7 / 0,1,4,5 & 1,7,7,3,9,3 \\
159 & 810|870|923|232|292|189 / 002110 & 144 / 3 & ES / 6 / A & 1,2,6,7 / 0,1,4,5 & 3,9,3,1,7,7 \\
160 & 890|850|963|252|212|149 / 210022 & 144 / 3 & ES / 6 / A & 1,2,6,7 / 0,1,4,5 & 1,7,7,3,9,3 \\
161 & 272|272|169|810|810|983 / 221002 & 144 / 3 & NO / 8 / A & 1,2,6,7 / 0,1,4,5 & 5,5,5,3,3,9 \\
162 & 252|252|109|830|830|943 / 001221 & 144 / 3 & NO / 8 / A & 1,2,6,7 / 0,1,4,5 & 3,3,9,5,5,5 \\
163 & 830|830|943|252|252|109 / 221001 & 144 / 3 & ES / 8 / A & 1,2,6,7 / 0,1,4,5 & 5,5,5,3,3,9 \\
164 & 810|810|983|272|272|169 / 002221 & 144 / 3 & ES / 8 / A & 1,2,6,7 / 0,1,4,5 & 3,3,9,5,5,5 \\
165 & 227|227|114|885|885|998 / 220002 & 1944 / 3 & NO / 45 / H & 0,1,2,3,4,5,6,7,8 / 0,1,2,3,4,5,6,7,8 & 0,0,0,0,0,0 \\
166 & 292|292|129|890|850|963 / 110210 & 144 / 3 & NO / 8 / B & 0,3,5,8 / 0,1,4,6 & 7,7,1,1,7,7 \\
167 & 232|292|189|850|850|903 / 110110 & 144 / 3 & NO / 8 / B & 0,3,5,8 / 0,1,4,6 & 1,7,7,7,7,1 \\
168 & 885|885|998|227|227|114 / 002220 & 1944 / 3 & ES / 45 / H & 0,1,2,3,4,5,6,7,8 / 0,1,2,3,4,5,6,7,8 & 0,0,0,0,0,0 \\
169 & 850|850|903|232|292|189 / 110110 & 144 / 3 & ES / 8 / B & 0,3,5,8 / 0,1,4,6 & 7,7,1,1,7,7 \\
170 & 890|850|963|292|292|129 / 210110 & 144 / 3 & ES / 8 / B & 0,3,5,8 / 0,1,4,6 & 1,7,7,7,7,1 \\
171 & 292|232|189|830|830|943 / 110221 & 144 / 3 & NO / 8 / A & 0,1,7,8 / 0,1,2,3 & 7,1,7,5,5,5 \\
172 & 292|292|129|810|810|983 / 110002 & 144 / 3 & NO / 6 / A & 0,4,8 / 0,1,5 & 7,7,1,3,3,9 \\
173 & 272|272|169|850|890|963 / 221120 & 144 / 3 & NO / 8 / A & 0,1,7,8 / 0,1,2,3 & 5,5,5,7,1,7 \\
174 & 252|252|109|850|850|903 / 001110 & 144 / 3 & NO / 6 / A & 0,4,8 / 0,1,5 & 3,3,9,7,7,1 \\
175 & 850|890|963|272|272|169 / 120221 & 144 / 3 & ES / 8 / A & 0,1,7,8 / 0,1,2,3 & 7,1,7,5,5,5 \\
176 & 850|850|903|252|252|109 / 110001 & 144 / 3 & ES / 6 / A & 0,4,8 / 0,1,5 & 7,7,1,3,3,9 \\
177 & 830|830|943|292|232|189 / 221110 & 144 / 3 & ES / 8 / A & 0,1,7,8 / 0,1,2,3 & 5,5,5,7,1,7 \\
178 & 810|810|983|292|292|129 / 002110 & 144 / 3 & ES / 6 / A & 0,4,8 / 0,1,5 & 3,3,9,7,7,1 \\
179 & 292|292|129|830|830|943 / 110221 & 144 / 3 & NO / 6 / A & 1,3,5,7 / 0,2,4,6 & 7,7,1,5,5,5 \\
180 & 212|252|149|850|890|963 / 202120 & 144 / 3 & NO / 8 / A & 1,3,5,7 / 0,2,4,6 & 9,3,3,7,1,7 \\
181 & 272|272|169|850|850|903 / 221110 & 144 / 3 & NO / 6 / A & 1,3,5,7 / 0,2,4,6 & 5,5,5,7,7,1 \\
182 & 292|232|189|870|810|923 / 110002 & 144 / 3 & NO / 8 / A & 1,3,5,7 / 0,2,4,6 & 7,1,7,9,3,3 \\
183 & 850|850|903|272|272|169 / 110221 & 144 / 3 & ES / 6 / A & 1,3,5,7 / 0,2,4,6 & 7,7,1,5,5,5 \\
184 & 870|810|923|292|232|189 / 002110 & 144 / 3 & ES / 8 / A & 1,3,5,7 / 0,2,4,6 & 9,3,3,7,1,7 \\
185 & 830|830|943|292|292|129 / 221110 & 144 / 3 & ES / 6 / A & 1,3,5,7 / 0,2,4,6 & 5,5,5,7,7,1 \\
186 & 850|890|963|212|252|149 / 120202 & 144 / 3 & ES / 8 / A & 1,3,5,7 / 0,2,4,6 & 7,1,7,9,3,3 \\
187 & 292|232|189|850|890|963 / 110120 & 144 / 3 & NO / 4 / B & 2,6 / 0,4 & 7,1,7,7,1,7 \\
188 & 850|890|963|292|232|189 / 120110 & 144 / 3 & ES / 4 / B & 2,6 / 0,4 & 7,1,7,7,1,7 \\
189 & 272|272|169|870|810|923 / 221002 & 144 / 3 & NO / 6 / A & 3,4,5 / 0,1,2 & 5,5,5,9,3,3 \\
190 & 212|252|149|830|830|943 / 202221 & 144 / 3 & NO / 6 / A & 3,4,5 / 0,1,2 & 9,3,3,5,5,5 \\
191 & 830|830|943|212|252|149 / 221202 & 144 / 3 & ES / 6 / A & 3,4,5 / 0,1,2 & 5,5,5,9,3,3 \\
192 & 870|810|923|272|272|169 / 002221 & 144 / 3 & ES / 6 / A & 3,4,5 / 0,1,2 & 9,3,3,5,5,5 \\
193 & 252|212|149|830|830|943 / 022221 & 144 / 3 & NO / 6 / A & 0,1,7,8 / 0,1,2,3 & 3,9,3,5,5,5 \\
194 & 272|272|169|810|870|923 / 221002 & 144 / 3 & NO / 6 / A & 0,1,7,8 / 0,1,2,3 & 5,5,5,3,9,3 \\
195 & 810|870|923|272|272|169 / 002221 & 144 / 3 & ES / 6 / A & 0,1,7,8 / 0,1,2,3 & 3,9,3,5,5,5 \\
196 & 830|830|943|252|212|149 / 221022 & 144 / 3 & ES / 6 / A & 0,1,7,8 / 0,1,2,3 & 5,5,5,3,9,3 \\
197 & 252|212|149|810|870|923 / 022002 & 144 / 3 & NO / 4 / B & 0,8 / 0,1 & 3,9,3,3,9,3 \\
198 & 810|870|923|252|212|149 / 002022 & 144 / 3 & ES / 4 / B & 0,8 / 0,1 & 3,9,3,3,9,3 \\
199 & 272|272|169|890|850|963 / 221210 & 144 / 3 & NO / 6 / A & 2,4,6 / 0,2,4 & 5,5,5,1,7,7 \\
200 & 232|292|189|830|830|943 / 110221 & 144 / 3 & NO / 6 / A & 2,4,6 / 0,2,4 & 1,7,7,5,5,5 \\
201 & 830|830|943|232|292|189 / 221110 & 144 / 3 & ES / 6 / A & 2,4,6 / 0,2,4 & 5,5,5,1,7,7 \\
202 & 890|850|963|272|272|169 / 210221 & 144 / 3 & ES / 6 / A & 2,4,6 / 0,2,4 & 1,7,7,5,5,5 \\
203 & 212|252|149|810|810|983 / 202002 & 144 / 3 & NO / 8 / B & 2,3,5,6 / 0,1,3,4 & 9,3,3,3,3,9 \\
204 & 252|252|109|870|810|923 / 001002 & 144 / 3 & NO / 8 / B & 2,3,5,6 / 0,1,3,4 & 3,3,9,9,3,3 \\
205 & 870|810|923|252|252|109 / 002001 & 144 / 3 & ES / 8 / B & 2,3,5,6 / 0,1,3,4 & 9,3,3,3,3,9 \\
206 & 810|810|983|212|252|149 / 002202 & 144 / 3 & ES / 8 / B & 2,3,5,6 / 0,1,3,4 & 3,3,9,9,3,3 \\
207 & 272|272|169|830|830|943 / 221221 & 432 / 3 & NO / 12 / B & 0,2,3,5,6,8 / 0,1,3,4,6,7 & 5,5,5,5,5,5 \\
208 & 830|830|943|272|272|169 / 221221 & 432 / 3 & ES / 12 / B & 0,2,3,5,6,8 / 0,1,3,4,6,7 & 5,5,5,5,5,5 \\
209 & 581|654|478|129|292|292 / 201011 & 144 / 4 & NO / 6 / A & 0,1,5,6 / 0,1,4,5 & 3,9,3,1,7,7 \\
210 & 561|634|418|149|212|252 / 011220 & 144 / 4 & NO / 6 / A & 0,1,5,6 / 0,1,4,5 & 1,7,7,3,9,3 \\
211 & 149|212|252|561|634|418 / 220011 & 144 / 4 & ES / 6 / A & 0,1,5,6 / 0,1,4,5 & 3,9,3,1,7,7 \\
212 & 129|292|292|581|654|478 / 011201 & 144 / 4 & ES / 6 / A & 0,1,5,6 / 0,1,4,5 & 1,7,7,3,9,3 \\
213 & 501|614|498|149|252|212 / 020202 & 144 / 4 & NO / 8 / A & 0,1,5,6 / 0,1,4,5 & 5,5,5,3,3,9 \\
214 & 581|694|438|169|272|272 / 210122 & 144 / 4 & NO / 8 / A & 0,1,5,6 / 0,1,4,5 & 3,3,9,5,5,5 \\
215 & 169|272|272|581|694|438 / 122210 & 144 / 4 & ES / 8 / A & 0,1,5,6 / 0,1,4,5 & 5,5,5,3,3,9 \\
216 & 149|252|212|501|614|498 / 202020 & 144 / 4 & ES / 8 / A & 0,1,5,6 / 0,1,4,5 & 3,3,9,5,5,5 \\
217 & 556|669|443|114|227|227 / 102022 & 1944 / 4 & NO / 45 / H & 0,1,2,3,4,5,6,7,8 / 0,1,2,3,4,5,6,7,8 & 0,0,0,0,0,0 \\
218 & 521|634|458|129|292|292 / 212011 & 144 / 4 & NO / 8 / B & 2,4,7,8 / 0,1,4,6 & 7,7,1,1,7,7 \\
219 & 561|634|418|189|292|232 / 011011 & 144 / 4 & NO / 8 / B & 2,4,7,8 / 0,1,4,6 & 1,7,7,7,7,1 \\
220 & 114|227|227|556|669|443 / 022102 & 1944 / 4 & ES / 45 / H & 0,1,2,3,4,5,6,7,8 / 0,1,2,3,4,5,6,7,8 & 0,0,0,0,0,0 \\
221 & 189|292|232|561|634|418 / 011011 & 144 / 4 & ES / 8 / B & 2,4,7,8 / 0,1,4,6 & 7,7,1,1,7,7 \\
222 & 129|292|292|521|634|458 / 011212 & 144 / 4 & ES / 8 / B & 2,4,7,8 / 0,1,4,6 & 1,7,7,7,7,1 \\
223 & 521|674|418|169|272|272 / 221122 & 144 / 4 & NO / 8 / A & 0,6,7,8 / 0,1,2,3 & 7,1,7,5,5,5 \\
224 & 521|634|458|149|252|212 / 212202 & 144 / 4 & NO / 6 / A & 3,7,8 / 0,1,5 & 7,7,1,3,3,9 \\
225 & 501|614|498|189|232|292 / 020011 & 144 / 4 & NO / 8 / A & 0,6,7,8 / 0,1,2,3 & 5,5,5,7,1,7 \\
226 & 581|694|438|189|292|232 / 210011 & 144 / 4 & NO / 6 / A & 3,7,8 / 0,1,5 & 3,3,9,7,7,1 \\
227 & 189|232|292|501|614|498 / 011020 & 144 / 4 & ES / 8 / A & 0,6,7,8 / 0,1,2,3 & 7,1,7,5,5,5 \\
228 & 189|292|232|581|694|438 / 011210 & 144 / 4 & ES / 6 / A & 3,7,8 / 0,1,5 & 7,7,1,3,3,9 \\
229 & 169|272|272|521|674|418 / 122221 & 144 / 4 & ES / 8 / A & 0,6,7,8 / 0,1,2,3 & 5,5,5,7,1,7 \\
230 & 149|252|212|521|634|458 / 202212 & 144 / 4 & ES / 6 / A & 3,7,8 / 0,1,5 & 3,3,9,7,7,1 \\
231 & 521|634|458|169|272|272 / 212122 & 144 / 4 & NO / 6 / A & 0,2,4,6 / 0,2,4,6 & 7,7,1,5,5,5 \\
232 & 541|694|478|189|232|292 / 111011 & 144 / 4 & NO / 8 / A & 0,2,4,6 / 0,2,4,6 & 9,3,3,7,1,7 \\
233 & 501|614|498|189|292|232 / 020011 & 144 / 4 & NO / 6 / A & 0,2,4,6 / 0,2,4,6 & 5,5,5,7,7,1 \\
234 & 521|674|418|109|252|252 / 221100 & 144 / 4 & NO / 8 / A & 0,2,4,6 / 0,2,4,6 & 7,1,7,9,3,3 \\
235 & 189|292|232|501|614|498 / 011020 & 144 / 4 & ES / 6 / A & 0,2,4,6 / 0,2,4,6 & 7,7,1,5,5,5 \\
236 & 109|252|252|521|674|418 / 100221 & 144 / 4 & ES / 8 / A & 0,2,4,6 / 0,2,4,6 & 9,3,3,7,1,7 \\
237 & 169|272|272|521|634|458 / 122212 & 144 / 4 & ES / 6 / A & 0,2,4,6 / 0,2,4,6 & 5,5,5,7,7,1 \\
238 & 189|232|292|541|694|478 / 011111 & 144 / 4 & ES / 8 / A & 0,2,4,6 / 0,2,4,6 & 7,1,7,9,3,3 \\
239 & 521|674|418|189|232|292 / 221011 & 144 / 4 & NO / 4 / B & 1,5 / 0,4 & 7,1,7,7,1,7 \\
240 & 189|232|292|521|674|418 / 011221 & 144 / 4 & ES / 4 / B & 1,5 / 0,4 & 7,1,7,7,1,7 \\
241 & 501|614|498|109|252|252 / 020100 & 144 / 4 & NO / 6 / A & 2,3,4 / 0,1,2 & 5,5,5,9,3,3 \\
242 & 541|694|478|169|272|272 / 111122 & 144 / 4 & NO / 6 / A & 2,3,4 / 0,1,2 & 9,3,3,5,5,5 \\
243 & 169|272|272|541|694|478 / 122111 & 144 / 4 & ES / 6 / A & 2,3,4 / 0,1,2 & 5,5,5,9,3,3 \\
244 & 109|252|252|501|614|498 / 100020 & 144 / 4 & ES / 6 / A & 2,3,4 / 0,1,2 & 9,3,3,5,5,5 \\
245 & 581|654|478|169|272|272 / 201122 & 144 / 4 & NO / 6 / A & 0,6,7,8 / 0,1,2,3 & 3,9,3,5,5,5 \\
246 & 501|614|498|149|212|252 / 020220 & 144 / 4 & NO / 6 / A & 0,6,7,8 / 0,1,2,3 & 5,5,5,3,9,3 \\
247 & 149|212|252|501|614|498 / 220020 & 144 / 4 & ES / 6 / A & 0,6,7,8 / 0,1,2,3 & 3,9,3,5,5,5 \\
248 & 169|272|272|581|654|478 / 122201 & 144 / 4 & ES / 6 / A & 0,6,7,8 / 0,1,2,3 & 5,5,5,3,9,3 \\
249 & 581|654|478|149|212|252 / 201220 & 144 / 4 & NO / 4 / B & 7,8 / 0,1 & 3,9,3,3,9,3 \\
250 & 149|212|252|581|654|478 / 220201 & 144 / 4 & ES / 4 / B & 7,8 / 0,1 & 3,9,3,3,9,3 \\
251 & 501|614|498|129|292|292 / 020011 & 144 / 4 & NO / 6 / A & 1,3,5 / 0,2,4 & 5,5,5,1,7,7 \\
252 & 561|634|418|169|272|272 / 011122 & 144 / 4 & NO / 6 / A & 1,3,5 / 0,2,4 & 1,7,7,5,5,5 \\
253 & 169|272|272|561|634|418 / 122011 & 144 / 4 & ES / 6 / A & 1,3,5 / 0,2,4 & 5,5,5,1,7,7 \\
254 & 129|292|292|501|614|498 / 011020 & 144 / 4 & ES / 6 / A & 1,3,5 / 0,2,4 & 1,7,7,5,5,5 \\
255 & 541|694|478|149|252|212 / 111202 & 144 / 4 & NO / 8 / B & 1,2,4,5 / 0,1,3,4 & 9,3,3,3,3,9 \\
256 & 581|694|438|109|252|252 / 210100 & 144 / 4 & NO / 8 / B & 1,2,4,5 / 0,1,3,4 & 3,3,9,9,3,3 \\
257 & 109|252|252|581|694|438 / 100210 & 144 / 4 & ES / 8 / B & 1,2,4,5 / 0,1,3,4 & 9,3,3,3,3,9 \\
258 & 149|252|212|541|694|478 / 202111 & 144 / 4 & ES / 8 / B & 1,2,4,5 / 0,1,3,4 & 3,3,9,9,3,3 \\
259 & 501|614|498|169|272|272 / 020122 & 432 / 4 & NO / 12 / B & 1,2,4,5,7,8 / 0,1,3,4,6,7 & 5,5,5,5,5,5 \\
260 & 169|272|272|501|614|498 / 122020 & 432 / 4 & ES / 12 / B & 1,2,4,5,7,8 / 0,1,3,4,6,7 & 5,5,5,5,5,5 \\
261 & 810|983|810|458|634|521 / 020212 & 144 / 5 & NO / 6 / A & 0,4,5,8 / 0,1,4,5 & 3,9,3,1,7,7 \\
262 & 890|963|850|478|654|581 / 201102 & 144 / 5 & NO / 6 / A & 0,4,5,8 / 0,1,4,5 & 1,7,7,3,9,3 \\
263 & 478|654|581|890|963|850 / 102201 & 144 / 5 & ES / 6 / A & 0,4,5,8 / 0,1,4,5 & 3,9,3,1,7,7 \\
264 & 458|634|521|810|983|810 / 212020 & 144 / 5 & ES / 6 / A & 0,4,5,8 / 0,1,4,5 & 1,7,7,3,9,3 \\
265 & 830|943|830|478|694|541 / 212111 & 144 / 5 & NO / 8 / A & 0,4,5,8 / 0,1,4,5 & 5,5,5,3,3,9 \\
266 & 810|923|870|498|614|501 / 020020 & 144 / 5 & NO / 8 / A & 0,4,5,8 / 0,1,4,5 & 3,3,9,5,5,5 \\
267 & 498|614|501|810|923|870 / 020020 & 144 / 5 & ES / 8 / A & 0,4,5,8 / 0,1,4,5 & 5,5,5,3,3,9 \\
268 & 478|694|541|830|943|830 / 111212 & 144 / 5 & ES / 8 / A & 0,4,5,8 / 0,1,4,5 & 3,3,9,5,5,5 \\
269 & 885|998|885|443|669|556 / 020201 & 1944 / 5 & NO / 45 / H & 0,1,2,3,4,5,6,7,8 / 0,1,2,3,4,5,6,7,8 & 0,0,0,0,0,0 \\
270 & 850|963|890|458|634|521 / 102212 & 144 / 5 & NO / 8 / B & 1,3,6,7 / 0,1,4,6 & 7,7,1,1,7,7 \\
271 & 890|963|850|418|634|561 / 201110 & 144 / 5 & NO / 8 / B & 1,3,6,7 / 0,1,4,6 & 1,7,7,7,7,1 \\
272 & 443|669|556|885|998|885 / 201020 & 1944 / 5 & ES / 45 / H & 0,1,2,3,4,5,6,7,8 / 0,1,2,3,4,5,6,7,8 & 0,0,0,0,0,0 \\
273 & 418|634|561|890|963|850 / 110201 & 144 / 5 & ES / 8 / B & 1,3,6,7 / 0,1,4,6 & 7,7,1,1,7,7 \\
274 & 458|634|521|850|963|890 / 212102 & 144 / 5 & ES / 8 / B & 1,3,6,7 / 0,1,4,6 & 1,7,7,7,7,1 \\
275 & 850|903|850|498|614|501 / 101020 & 144 / 5 & NO / 8 / A & 5,6,7,8 / 0,1,2,3 & 7,1,7,5,5,5 \\
276 & 850|963|890|478|694|541 / 102111 & 144 / 5 & NO / 6 / A & 2,6,7 / 0,1,5 & 7,7,1,3,3,9 \\
277 & 830|943|830|418|674|521 / 212122 & 144 / 5 & NO / 8 / A & 5,6,7,8 / 0,1,2,3 & 5,5,5,7,1,7 \\
278 & 810|923|870|418|634|561 / 020110 & 144 / 5 & NO / 6 / A & 2,6,7 / 0,1,5 & 3,3,9,7,7,1 \\
279 & 418|674|521|830|943|830 / 122212 & 144 / 5 & ES / 8 / A & 5,6,7,8 / 0,1,2,3 & 7,1,7,5,5,5 \\
280 & 418|634|561|810|923|870 / 110020 & 144 / 5 & ES / 6 / A & 2,6,7 / 0,1,5 & 7,7,1,3,3,9 \\
281 & 498|614|501|850|903|850 / 020101 & 144 / 5 & ES / 8 / A & 5,6,7,8 / 0,1,2,3 & 5,5,5,7,1,7 \\
282 & 478|694|541|850|963|890 / 111102 & 144 / 5 & ES / 6 / A & 2,6,7 / 0,1,5 & 3,3,9,7,7,1 \\
283 & 850|963|890|498|614|501 / 102020 & 144 / 5 & NO / 6 / A & 1,3,5,8 / 0,2,4,6 & 7,7,1,5,5,5 \\
284 & 870|923|810|418|674|521 / 020122 & 144 / 5 & NO / 8 / A & 1,3,5,8 / 0,2,4,6 & 9,3,3,7,1,7 \\
285 & 830|943|830|418|634|561 / 212110 & 144 / 5 & NO / 6 / A & 1,3,5,8 / 0,2,4,6 & 5,5,5,7,7,1 \\
286 & 850|903|850|438|694|581 / 101012 & 144 / 5 & NO / 8 / A & 1,3,5,8 / 0,2,4,6 & 7,1,7,9,3,3 \\
287 & 418|634|561|830|943|830 / 110212 & 144 / 5 & ES / 6 / A & 1,3,5,8 / 0,2,4,6 & 7,7,1,5,5,5 \\
288 & 438|694|581|850|903|850 / 012101 & 144 / 5 & ES / 8 / A & 1,3,5,8 / 0,2,4,6 & 9,3,3,7,1,7 \\
289 & 498|614|501|850|963|890 / 020102 & 144 / 5 & ES / 6 / A & 1,3,5,8 / 0,2,4,6 & 5,5,5,7,7,1 \\
290 & 418|674|521|870|923|810 / 122020 & 144 / 5 & ES / 8 / A & 1,3,5,8 / 0,2,4,6 & 7,1,7,9,3,3 \\
291 & 850|903|850|418|674|521 / 101122 & 144 / 5 & NO / 4 / B & 0,4 / 0,4 & 7,1,7,7,1,7 \\
292 & 418|674|521|850|903|850 / 122101 & 144 / 5 & ES / 4 / B & 0,4 / 0,4 & 7,1,7,7,1,7 \\
293 & 830|943|830|438|694|581 / 212012 & 144 / 5 & NO / 6 / A & 1,2,3 / 0,1,2 & 5,5,5,9,3,3 \\
294 & 870|923|810|498|614|501 / 020020 & 144 / 5 & NO / 6 / A & 1,2,3 / 0,1,2 & 9,3,3,5,5,5 \\
295 & 498|614|501|870|923|810 / 020020 & 144 / 5 & ES / 6 / A & 1,2,3 / 0,1,2 & 5,5,5,9,3,3 \\
296 & 438|694|581|830|943|830 / 012212 & 144 / 5 & ES / 6 / A & 1,2,3 / 0,1,2 & 9,3,3,5,5,5 \\
297 & 810|983|810|498|614|501 / 020020 & 144 / 5 & NO / 6 / A & 5,6,7,8 / 0,1,2,3 & 3,9,3,5,5,5 \\
298 & 830|943|830|478|654|581 / 212102 & 144 / 5 & NO / 6 / A & 5,6,7,8 / 0,1,2,3 & 5,5,5,3,9,3 \\
299 & 478|654|581|830|943|830 / 102212 & 144 / 5 & ES / 6 / A & 5,6,7,8 / 0,1,2,3 & 3,9,3,5,5,5 \\
300 & 498|614|501|810|983|810 / 020020 & 144 / 5 & ES / 6 / A & 5,6,7,8 / 0,1,2,3 & 5,5,5,3,9,3 \\
301 & 810|983|810|478|654|581 / 020102 & 144 / 5 & NO / 4 / B & 6,7 / 0,1 & 3,9,3,3,9,3 \\
302 & 478|654|581|810|983|810 / 102020 & 144 / 5 & ES / 4 / B & 6,7 / 0,1 & 3,9,3,3,9,3 \\
303 & 830|943|830|458|634|521 / 212212 & 144 / 5 & NO / 6 / A & 0,2,4 / 0,2,4 & 5,5,5,1,7,7 \\
304 & 890|963|850|498|614|501 / 201020 & 144 / 5 & NO / 6 / A & 0,2,4 / 0,2,4 & 1,7,7,5,5,5 \\
305 & 498|614|501|890|963|850 / 020201 & 144 / 5 & ES / 6 / A & 0,2,4 / 0,2,4 & 5,5,5,1,7,7 \\
306 & 458|634|521|830|943|830 / 212212 & 144 / 5 & ES / 6 / A & 0,2,4 / 0,2,4 & 1,7,7,5,5,5 \\
307 & 870|923|810|478|694|541 / 020111 & 144 / 5 & NO / 8 / B & 0,1,3,4 / 0,1,3,4 & 9,3,3,3,3,9 \\
308 & 810|923|870|438|694|581 / 020012 & 144 / 5 & NO / 8 / B & 0,1,3,4 / 0,1,3,4 & 3,3,9,9,3,3 \\
309 & 438|694|581|810|923|870 / 012020 & 144 / 5 & ES / 8 / B & 0,1,3,4 / 0,1,3,4 & 9,3,3,3,3,9 \\
310 & 478|694|541|870|923|810 / 111020 & 144 / 5 & ES / 8 / B & 0,1,3,4 / 0,1,3,4 & 3,3,9,9,3,3 \\
311 & 830|943|830|498|614|501 / 212020 & 432 / 5 & NO / 12 / B & 0,1,3,4,6,7 / 0,1,3,4,6,7 & 5,5,5,5,5,5 \\
312 & 498|614|501|830|943|830 / 020212 & 432 / 5 & ES / 12 / B & 0,1,3,4,6,7 / 0,1,3,4,6,7 & 5,5,5,5,5,5 \\
313 & 149|212|252|890|963|850 / 220201 & 144 / 6 & NO / 6 / A & 3,4,7,8 / 0,1,4,5 & 3,9,3,1,7,7 \\
314 & 129|292|292|810|983|810 / 011020 & 144 / 6 & NO / 6 / A & 3,4,7,8 / 0,1,4,5 & 1,7,7,3,9,3 \\
315 & 810|983|810|129|292|292 / 020011 & 144 / 6 & ES / 6 / A & 3,4,7,8 / 0,1,4,5 & 3,9,3,1,7,7 \\
316 & 890|963|850|149|212|252 / 201220 & 144 / 6 & ES / 6 / A & 3,4,7,8 / 0,1,4,5 & 1,7,7,3,9,3 \\
317 & 169|272|272|810|923|870 / 122020 & 144 / 6 & NO / 8 / A & 3,4,7,8 / 0,1,4,5 & 5,5,5,3,3,9 \\
318 & 149|252|212|830|943|830 / 202212 & 144 / 6 & NO / 8 / A & 3,4,7,8 / 0,1,4,5 & 3,3,9,5,5,5 \\
319 & 830|943|830|149|252|212 / 212202 & 144 / 6 & ES / 8 / A & 3,4,7,8 / 0,1,4,5 & 5,5,5,3,3,9 \\
320 & 810|923|870|169|272|272 / 020122 & 144 / 6 & ES / 8 / A & 3,4,7,8 / 0,1,4,5 & 3,3,9,5,5,5 \\
321 & 114|227|227|885|998|885 / 022020 & 1944 / 6 & NO / 45 / H & 0,1,2,3,4,5,6,7,8 / 0,1,2,3,4,5,6,7,8 & 0,0,0,0,0,0 \\
322 & 189|292|232|890|963|850 / 011201 & 144 / 6 & NO / 8 / B & 0,2,5,6 / 0,1,4,6 & 7,7,1,1,7,7 \\
323 & 129|292|292|850|963|890 / 011102 & 144 / 6 & NO / 8 / B & 0,2,5,6 / 0,1,4,6 & 1,7,7,7,7,1 \\
324 & 885|998|885|114|227|227 / 020022 & 1944 / 6 & ES / 45 / H & 0,1,2,3,4,5,6,7,8 / 0,1,2,3,4,5,6,7,8 & 0,0,0,0,0,0 \\
325 & 850|963|890|129|292|292 / 102011 & 144 / 6 & ES / 8 / B & 0,2,5,6 / 0,1,4,6 & 7,7,1,1,7,7 \\
326 & 890|963|850|189|292|232 / 201011 & 144 / 6 & ES / 8 / B & 0,2,5,6 / 0,1,4,6 & 1,7,7,7,7,1 \\
327 & 189|232|292|830|943|830 / 011212 & 144 / 6 & NO / 8 / A & 4,5,6,7 / 0,1,2,3 & 7,1,7,5,5,5 \\
328 & 189|292|232|810|923|870 / 011020 & 144 / 6 & NO / 6 / A & 1,5,6 / 0,1,5 & 7,7,1,3,3,9 \\
329 & 169|272|272|850|903|850 / 122101 & 144 / 6 & NO / 8 / A & 4,5,6,7 / 0,1,2,3 & 5,5,5,7,1,7 \\
330 & 149|252|212|850|963|890 / 202102 & 144 / 6 & NO / 6 / A & 1,5,6 / 0,1,5 & 3,3,9,7,7,1 \\
331 & 850|903|850|169|272|272 / 101122 & 144 / 6 & ES / 8 / A & 4,5,6,7 / 0,1,2,3 & 7,1,7,5,5,5 \\
332 & 850|963|890|149|252|212 / 102202 & 144 / 6 & ES / 6 / A & 1,5,6 / 0,1,5 & 7,7,1,3,3,9 \\
333 & 830|943|830|189|232|292 / 212011 & 144 / 6 & ES / 8 / A & 4,5,6,7 / 0,1,2,3 & 5,5,5,7,1,7 \\
334 & 810|923|870|189|292|232 / 020011 & 144 / 6 & ES / 6 / A & 1,5,6 / 0,1,5 & 3,3,9,7,7,1 \\
335 & 189|292|232|830|943|830 / 011212 & 144 / 6 & NO / 6 / A & 0,2,4,7 / 0,2,4,6 & 7,7,1,5,5,5 \\
336 & 109|252|252|850|903|850 / 100101 & 144 / 6 & NO / 8 / A & 0,2,4,7 / 0,2,4,6 & 9,3,3,7,1,7 \\
337 & 169|272|272|850|963|890 / 122102 & 144 / 6 & NO / 6 / A & 0,2,4,7 / 0,2,4,6 & 5,5,5,7,7,1 \\
338 & 189|232|292|870|923|810 / 011020 & 144 / 6 & NO / 8 / A & 0,2,4,7 / 0,2,4,6 & 7,1,7,9,3,3 \\
339 & 850|963|890|169|272|272 / 102122 & 144 / 6 & ES / 6 / A & 0,2,4,7 / 0,2,4,6 & 7,7,1,5,5,5 \\
340 & 870|923|810|189|232|292 / 020011 & 144 / 6 & ES / 8 / A & 0,2,4,7 / 0,2,4,6 & 9,3,3,7,1,7 \\
341 & 830|943|830|189|292|232 / 212011 & 144 / 6 & ES / 6 / A & 0,2,4,7 / 0,2,4,6 & 5,5,5,7,7,1 \\
342 & 850|903|850|109|252|252 / 101100 & 144 / 6 & ES / 8 / A & 0,2,4,7 / 0,2,4,6 & 7,1,7,9,3,3 \\
343 & 189|232|292|850|903|850 / 011101 & 144 / 6 & NO / 4 / B & 3,8 / 0,4 & 7,1,7,7,1,7 \\
344 & 850|903|850|189|232|292 / 101011 & 144 / 6 & ES / 4 / B & 3,8 / 0,4 & 7,1,7,7,1,7 \\
345 & 169|272|272|870|923|810 / 122020 & 144 / 6 & NO / 6 / A & 0,1,2 / 0,1,2 & 5,5,5,9,3,3 \\
346 & 109|252|252|830|943|830 / 100212 & 144 / 6 & NO / 6 / A & 0,1,2 / 0,1,2 & 9,3,3,5,5,5 \\
347 & 830|943|830|109|252|252 / 212100 & 144 / 6 & ES / 6 / A & 0,1,2 / 0,1,2 & 5,5,5,9,3,3 \\
348 & 870|923|810|169|272|272 / 020122 & 144 / 6 & ES / 6 / A & 0,1,2 / 0,1,2 & 9,3,3,5,5,5 \\
349 & 149|212|252|830|943|830 / 220212 & 144 / 6 & NO / 6 / A & 4,5,6,7 / 0,1,2,3 & 3,9,3,5,5,5 \\
350 & 169|272|272|810|983|810 / 122020 & 144 / 6 & NO / 6 / A & 4,5,6,7 / 0,1,2,3 & 5,5,5,3,9,3 \\
351 & 810|983|810|169|272|272 / 020122 & 144 / 6 & ES / 6 / A & 4,5,6,7 / 0,1,2,3 & 3,9,3,5,5,5 \\
352 & 830|943|830|149|212|252 / 212220 & 144 / 6 & ES / 6 / A & 4,5,6,7 / 0,1,2,3 & 5,5,5,3,9,3 \\
353 & 149|212|252|810|983|810 / 220020 & 144 / 6 & NO / 4 / B & 5,6 / 0,1 & 3,9,3,3,9,3 \\
354 & 810|983|810|149|212|252 / 020220 & 144 / 6 & ES / 4 / B & 5,6 / 0,1 & 3,9,3,3,9,3 \\
355 & 169|272|272|890|963|850 / 122201 & 144 / 6 & NO / 6 / A & 1,3,8 / 0,2,4 & 5,5,5,1,7,7 \\
356 & 129|292|292|830|943|830 / 011212 & 144 / 6 & NO / 6 / A & 1,3,8 / 0,2,4 & 1,7,7,5,5,5 \\
357 & 830|943|830|129|292|292 / 212011 & 144 / 6 & ES / 6 / A & 1,3,8 / 0,2,4 & 5,5,5,1,7,7 \\
358 & 890|963|850|169|272|272 / 201122 & 144 / 6 & ES / 6 / A & 1,3,8 / 0,2,4 & 1,7,7,5,5,5 \\
359 & 109|252|252|810|923|870 / 100020 & 144 / 6 & NO / 8 / B & 0,2,3,8 / 0,1,3,4 & 9,3,3,3,3,9 \\
360 & 149|252|212|870|923|810 / 202020 & 144 / 6 & NO / 8 / B & 0,2,3,8 / 0,1,3,4 & 3,3,9,9,3,3 \\
361 & 870|923|810|149|252|212 / 020202 & 144 / 6 & ES / 8 / B & 0,2,3,8 / 0,1,3,4 & 9,3,3,3,3,9 \\
362 & 810|923|870|109|252|252 / 020100 & 144 / 6 & ES / 8 / B & 0,2,3,8 / 0,1,3,4 & 3,3,9,9,3,3 \\
363 & 169|272|272|830|943|830 / 122212 & 432 / 6 & NO / 12 / B & 0,2,3,5,6,8 / 0,1,3,4,6,7 & 5,5,5,5,5,5 \\
364 & 830|943|830|169|272|272 / 212122 & 432 / 6 & ES / 12 / B & 0,2,3,5,6,8 / 0,1,3,4,6,7 & 5,5,5,5,5,5 \\
365 & 478|541|694|232|292|189 / 111110 & 144 / 7 & NO / 6 / A & 2,3,6,7 / 0,1,4,5 & 3,9,3,1,7,7 \\
366 & 458|521|634|252|212|149 / 221022 & 144 / 7 & NO / 6 / A & 2,3,6,7 / 0,1,4,5 & 1,7,7,3,9,3 \\
367 & 252|212|149|458|521|634 / 022221 & 144 / 7 & ES / 6 / A & 2,3,6,7 / 0,1,4,5 & 3,9,3,1,7,7 \\
368 & 232|292|189|478|541|694 / 110111 & 144 / 7 & ES / 6 / A & 2,3,6,7 / 0,1,4,5 & 1,7,7,3,9,3 \\
369 & 498|501|614|252|252|109 / 002001 & 144 / 7 & NO / 8 / A & 2,3,6,7 / 0,1,4,5 & 5,5,5,3,3,9 \\
370 & 478|581|654|272|272|169 / 120221 & 144 / 7 & NO / 8 / A & 2,3,6,7 / 0,1,4,5 & 3,3,9,5,5,5 \\
371 & 272|272|169|478|581|654 / 221120 & 144 / 7 & ES / 8 / A & 2,3,6,7 / 0,1,4,5 & 5,5,5,3,3,9 \\
372 & 252|252|109|498|501|614 / 001002 & 144 / 7 & ES / 8 / A & 2,3,6,7 / 0,1,4,5 & 3,3,9,5,5,5 \\
373 & 443|556|669|227|227|114 / 210220 & 1944 / 7 & NO / 45 / H & 0,1,2,3,4,5,6,7,8 / 0,1,2,3,4,5,6,7,8 & 0,0,0,0,0,0 \\
374 & 418|521|674|232|292|189 / 122110 & 144 / 7 & NO / 8 / B & 1,4,5,8 / 0,1,4,6 & 7,7,1,1,7,7 \\
375 & 458|521|634|292|292|129 / 221110 & 144 / 7 & NO / 8 / B & 1,4,5,8 / 0,1,4,6 & 1,7,7,7,7,1 \\
376 & 227|227|114|443|556|669 / 220210 & 1944 / 7 & ES / 45 / H & 0,1,2,3,4,5,6,7,8 / 0,1,2,3,4,5,6,7,8 & 0,0,0,0,0,0 \\
377 & 292|292|129|458|521|634 / 110221 & 144 / 7 & ES / 8 / B & 1,4,5,8 / 0,1,4,6 & 7,7,1,1,7,7 \\
378 & 232|292|189|418|521|674 / 110122 & 144 / 7 & ES / 8 / B & 1,4,5,8 / 0,1,4,6 & 1,7,7,7,7,1 \\
379 & 418|561|634|272|272|169 / 101221 & 144 / 7 & NO / 8 / A & 3,4,5,6 / 0,1,2,3 & 7,1,7,5,5,5 \\
380 & 418|521|674|252|252|109 / 122001 & 144 / 7 & NO / 6 / A & 0,4,5 / 0,1,5 & 7,7,1,3,3,9 \\
381 & 498|501|614|292|232|189 / 002110 & 144 / 7 & NO / 8 / A & 3,4,5,6 / 0,1,2,3 & 5,5,5,7,1,7 \\
382 & 478|581|654|292|292|129 / 120110 & 144 / 7 & NO / 6 / A & 0,4,5 / 0,1,5 & 3,3,9,7,7,1 \\
383 & 292|232|189|498|501|614 / 110002 & 144 / 7 & ES / 8 / A & 3,4,5,6 / 0,1,2,3 & 7,1,7,5,5,5 \\
384 & 292|292|129|478|581|654 / 110120 & 144 / 7 & ES / 6 / A & 0,4,5 / 0,1,5 & 7,7,1,3,3,9 \\
385 & 272|272|169|418|561|634 / 221101 & 144 / 7 & ES / 8 / A & 3,4,5,6 / 0,1,2,3 & 5,5,5,7,1,7 \\
386 & 252|252|109|418|521|674 / 001122 & 144 / 7 & ES / 6 / A & 0,4,5 / 0,1,5 & 3,3,9,7,7,1 \\
387 & 418|521|674|272|272|169 / 122221 & 144 / 7 & NO / 6 / A & 1,3,6,8 / 0,2,4,6 & 7,7,1,5,5,5 \\
388 & 438|581|694|292|232|189 / 021110 & 144 / 7 & NO / 8 / A & 1,3,6,8 / 0,2,4,6 & 9,3,3,7,1,7 \\
389 & 498|501|614|292|292|129 / 002110 & 144 / 7 & NO / 6 / A & 1,3,6,8 / 0,2,4,6 & 5,5,5,7,7,1 \\
390 & 418|561|634|212|252|149 / 101202 & 144 / 7 & NO / 8 / A & 1,3,6,8 / 0,2,4,6 & 7,1,7,9,3,3 \\
391 & 292|292|129|498|501|614 / 110002 & 144 / 7 & ES / 6 / A & 1,3,6,8 / 0,2,4,6 & 7,7,1,5,5,5 \\
392 & 212|252|149|418|561|634 / 202101 & 144 / 7 & ES / 8 / A & 1,3,6,8 / 0,2,4,6 & 9,3,3,7,1,7 \\
393 & 272|272|169|418|521|674 / 221122 & 144 / 7 & ES / 6 / A & 1,3,6,8 / 0,2,4,6 & 5,5,5,7,7,1 \\
394 & 292|232|189|438|581|694 / 110021 & 144 / 7 & ES / 8 / A & 1,3,6,8 / 0,2,4,6 & 7,1,7,9,3,3 \\
395 & 418|561|634|292|232|189 / 101110 & 144 / 7 & NO / 4 / B & 2,7 / 0,4 & 7,1,7,7,1,7 \\
396 & 292|232|189|418|561|634 / 110101 & 144 / 7 & ES / 4 / B & 2,7 / 0,4 & 7,1,7,7,1,7 \\
397 & 498|501|614|212|252|149 / 002202 & 144 / 7 & NO / 6 / A & 0,1,8 / 0,1,2 & 5,5,5,9,3,3 \\
398 & 438|581|694|272|272|169 / 021221 & 144 / 7 & NO / 6 / A & 0,1,8 / 0,1,2 & 9,3,3,5,5,5 \\
399 & 272|272|169|438|581|694 / 221021 & 144 / 7 & ES / 6 / A & 0,1,8 / 0,1,2 & 5,5,5,9,3,3 \\
400 & 212|252|149|498|501|614 / 202002 & 144 / 7 & ES / 6 / A & 0,1,8 / 0,1,2 & 9,3,3,5,5,5 \\
401 & 478|541|694|272|272|169 / 111221 & 144 / 7 & NO / 6 / A & 3,4,5,6 / 0,1,2,3 & 3,9,3,5,5,5 \\
402 & 498|501|614|252|212|149 / 002022 & 144 / 7 & NO / 6 / A & 3,4,5,6 / 0,1,2,3 & 5,5,5,3,9,3 \\
403 & 252|212|149|498|501|614 / 022002 & 144 / 7 & ES / 6 / A & 3,4,5,6 / 0,1,2,3 & 3,9,3,5,5,5 \\
404 & 272|272|169|478|541|694 / 221111 & 144 / 7 & ES / 6 / A & 3,4,5,6 / 0,1,2,3 & 5,5,5,3,9,3 \\
405 & 478|541|694|252|212|149 / 111022 & 144 / 7 & NO / 4 / B & 4,5 / 0,1 & 3,9,3,3,9,3 \\
406 & 252|212|149|478|541|694 / 022111 & 144 / 7 & ES / 4 / B & 4,5 / 0,1 & 3,9,3,3,9,3 \\
407 & 498|501|614|232|292|189 / 002110 & 144 / 7 & NO / 6 / A & 0,2,7 / 0,2,4 & 5,5,5,1,7,7 \\
408 & 458|521|634|272|272|169 / 221221 & 144 / 7 & NO / 6 / A & 0,2,7 / 0,2,4 & 1,7,7,5,5,5 \\
409 & 272|272|169|458|521|634 / 221221 & 144 / 7 & ES / 6 / A & 0,2,7 / 0,2,4 & 5,5,5,1,7,7 \\
410 & 232|292|189|498|501|614 / 110002 & 144 / 7 & ES / 6 / A & 0,2,7 / 0,2,4 & 1,7,7,5,5,5 \\
411 & 438|581|694|252|252|109 / 021001 & 144 / 7 & NO / 8 / B & 1,2,7,8 / 0,1,3,4 & 9,3,3,3,3,9 \\
412 & 478|581|654|212|252|149 / 120202 & 144 / 7 & NO / 8 / B & 1,2,7,8 / 0,1,3,4 & 3,3,9,9,3,3 \\
413 & 212|252|149|478|581|654 / 202120 & 144 / 7 & ES / 8 / B & 1,2,7,8 / 0,1,3,4 & 9,3,3,3,3,9 \\
414 & 252|252|109|438|581|694 / 001021 & 144 / 7 & ES / 8 / B & 1,2,7,8 / 0,1,3,4 & 3,3,9,9,3,3 \\
415 & 498|501|614|272|272|169 / 002221 & 432 / 7 & NO / 12 / B & 1,2,4,5,7,8 / 0,1,3,4,6,7 & 5,5,5,5,5,5 \\
416 & 272|272|169|498|501|614 / 221002 & 432 / 7 & ES / 12 / B & 1,2,4,5,7,8 / 0,1,3,4,6,7 & 5,5,5,5,5,5 \\
417 & 810|870|923|674|521|418 / 002221 & 144 / 8 & NO / 6 / A & 1,2,5,6 / 0,1,4,5 & 3,9,3,1,7,7 \\
418 & 890|850|963|694|541|478 / 210111 & 144 / 8 & NO / 6 / A & 1,2,5,6 / 0,1,4,5 & 1,7,7,3,9,3 \\
419 & 694|541|478|890|850|963 / 111210 & 144 / 8 & ES / 6 / A & 1,2,5,6 / 0,1,4,5 & 3,9,3,1,7,7 \\
420 & 674|521|418|810|870|923 / 221002 & 144 / 8 & ES / 6 / A & 1,2,5,6 / 0,1,4,5 & 1,7,7,3,9,3 \\
421 & 830|830|943|694|581|438 / 221120 & 144 / 8 & NO / 8 / A & 1,2,5,6 / 0,1,4,5 & 5,5,5,3,3,9 \\
422 & 810|810|983|614|501|498 / 002200 & 144 / 8 & NO / 8 / A & 1,2,5,6 / 0,1,4,5 & 3,3,9,5,5,5 \\
423 & 614|501|498|810|810|983 / 200002 & 144 / 8 & ES / 8 / A & 1,2,5,6 / 0,1,4,5 & 5,5,5,3,3,9 \\
424 & 694|581|438|830|830|943 / 120221 & 144 / 8 & ES / 8 / A & 1,2,5,6 / 0,1,4,5 & 3,3,9,5,5,5 \\
425 & 885|885|998|669|556|443 / 002012 & 1944 / 8 & NO / 45 / H & 0,1,2,3,4,5,6,7,8 / 0,1,2,3,4,5,6,7,8 & 0,0,0,0,0,0 \\
426 & 850|850|903|674|521|418 / 110221 & 144 / 8 & NO / 8 / B & 0,3,4,7 / 0,1,4,6 & 7,7,1,1,7,7 \\
427 & 890|850|963|634|521|458 / 210122 & 144 / 8 & NO / 8 / B & 0,3,4,7 / 0,1,4,6 & 1,7,7,7,7,1 \\
428 & 669|556|443|885|885|998 / 012002 & 1944 / 8 & ES / 45 / H & 0,1,2,3,4,5,6,7,8 / 0,1,2,3,4,5,6,7,8 & 0,0,0,0,0,0 \\
429 & 634|521|458|890|850|963 / 122210 & 144 / 8 & ES / 8 / B & 0,3,4,7 / 0,1,4,6 & 7,7,1,1,7,7 \\
430 & 674|521|418|850|850|903 / 221110 & 144 / 8 & ES / 8 / B & 0,3,4,7 / 0,1,4,6 & 1,7,7,7,7,1 \\
431 & 850|890|963|614|501|498 / 120200 & 144 / 8 & NO / 8 / A & 2,3,4,5 / 0,1,2,3 & 7,1,7,5,5,5 \\
432 & 850|850|903|694|581|438 / 110120 & 144 / 8 & NO / 6 / A & 3,4,8 / 0,1,5 & 7,7,1,3,3,9 \\
433 & 830|830|943|634|561|418 / 221101 & 144 / 8 & NO / 8 / A & 2,3,4,5 / 0,1,2,3 & 5,5,5,7,1,7 \\
434 & 810|810|983|634|521|458 / 002122 & 144 / 8 & NO / 6 / A & 3,4,8 / 0,1,5 & 3,3,9,7,7,1 \\
435 & 634|561|418|830|830|943 / 101221 & 144 / 8 & ES / 8 / A & 2,3,4,5 / 0,1,2,3 & 7,1,7,5,5,5 \\
436 & 634|521|458|810|810|983 / 122002 & 144 / 8 & ES / 6 / A & 3,4,8 / 0,1,5 & 7,7,1,3,3,9 \\
437 & 614|501|498|850|890|963 / 200120 & 144 / 8 & ES / 8 / A & 2,3,4,5 / 0,1,2,3 & 5,5,5,7,1,7 \\
438 & 694|581|438|850|850|903 / 120110 & 144 / 8 & ES / 6 / A & 3,4,8 / 0,1,5 & 3,3,9,7,7,1 \\
439 & 850|850|903|614|501|498 / 110200 & 144 / 8 & NO / 6 / A & 0,2,5,7 / 0,2,4,6 & 7,7,1,5,5,5 \\
440 & 870|810|923|634|561|418 / 002101 & 144 / 8 & NO / 8 / A & 0,2,5,7 / 0,2,4,6 & 9,3,3,7,1,7 \\
441 & 830|830|943|634|521|458 / 221122 & 144 / 8 & NO / 6 / A & 0,2,5,7 / 0,2,4,6 & 5,5,5,7,7,1 \\
442 & 850|890|963|654|581|478 / 120021 & 144 / 8 & NO / 8 / A & 0,2,5,7 / 0,2,4,6 & 7,1,7,9,3,3 \\
443 & 634|521|458|830|830|943 / 122221 & 144 / 8 & ES / 6 / A & 0,2,5,7 / 0,2,4,6 & 7,7,1,5,5,5 \\
444 & 654|581|478|850|890|963 / 021120 & 144 / 8 & ES / 8 / A & 0,2,5,7 / 0,2,4,6 & 9,3,3,7,1,7 \\
445 & 614|501|498|850|850|903 / 200110 & 144 / 8 & ES / 6 / A & 0,2,5,7 / 0,2,4,6 & 5,5,5,7,7,1 \\
446 & 634|561|418|870|810|923 / 101002 & 144 / 8 & ES / 8 / A & 0,2,5,7 / 0,2,4,6 & 7,1,7,9,3,3 \\
447 & 850|890|963|634|561|418 / 120101 & 144 / 8 & NO / 4 / B & 1,6 / 0,4 & 7,1,7,7,1,7 \\
448 & 634|561|418|850|890|963 / 101120 & 144 / 8 & ES / 4 / B & 1,6 / 0,4 & 7,1,7,7,1,7 \\
449 & 830|830|943|654|581|478 / 221021 & 144 / 8 & NO / 6 / A & 0,7,8 / 0,1,2 & 5,5,5,9,3,3 \\
450 & 870|810|923|614|501|498 / 002200 & 144 / 8 & NO / 6 / A & 0,7,8 / 0,1,2 & 9,3,3,5,5,5 \\
451 & 614|501|498|870|810|923 / 200002 & 144 / 8 & ES / 6 / A & 0,7,8 / 0,1,2 & 5,5,5,9,3,3 \\
452 & 654|581|478|830|830|943 / 021221 & 144 / 8 & ES / 6 / A & 0,7,8 / 0,1,2 & 9,3,3,5,5,5 \\
453 & 810|870|923|614|501|498 / 002200 & 144 / 8 & NO / 6 / A & 2,3,4,5 / 0,1,2,3 & 3,9,3,5,5,5 \\
454 & 830|830|943|694|541|478 / 221111 & 144 / 8 & NO / 6 / A & 2,3,4,5 / 0,1,2,3 & 5,5,5,3,9,3 \\
455 & 694|541|478|830|830|943 / 111221 & 144 / 8 & ES / 6 / A & 2,3,4,5 / 0,1,2,3 & 3,9,3,5,5,5 \\
456 & 614|501|498|810|870|923 / 200002 & 144 / 8 & ES / 6 / A & 2,3,4,5 / 0,1,2,3 & 5,5,5,3,9,3 \\
457 & 810|870|923|694|541|478 / 002111 & 144 / 8 & NO / 4 / B & 3,4 / 0,1 & 3,9,3,3,9,3 \\
458 & 694|541|478|810|870|923 / 111002 & 144 / 8 & ES / 4 / B & 3,4 / 0,1 & 3,9,3,3,9,3 \\
459 & 830|830|943|674|521|418 / 221221 & 144 / 8 & NO / 6 / A & 1,6,8 / 0,2,4 & 5,5,5,1,7,7 \\
460 & 890|850|963|614|501|498 / 210200 & 144 / 8 & NO / 6 / A & 1,6,8 / 0,2,4 & 1,7,7,5,5,5 \\
461 & 614|501|498|890|850|963 / 200210 & 144 / 8 & ES / 6 / A & 1,6,8 / 0,2,4 & 5,5,5,1,7,7 \\
462 & 674|521|418|830|830|943 / 221221 & 144 / 8 & ES / 6 / A & 1,6,8 / 0,2,4 & 1,7,7,5,5,5 \\
463 & 870|810|923|694|581|438 / 002120 & 144 / 8 & NO / 8 / B & 0,1,6,7 / 0,1,3,4 & 9,3,3,3,3,9 \\
464 & 810|810|983|654|581|478 / 002021 & 144 / 8 & NO / 8 / B & 0,1,6,7 / 0,1,3,4 & 3,3,9,9,3,3 \\
465 & 654|581|478|810|810|983 / 021002 & 144 / 8 & ES / 8 / B & 0,1,6,7 / 0,1,3,4 & 9,3,3,3,3,9 \\
466 & 694|581|438|870|810|923 / 120002 & 144 / 8 & ES / 8 / B & 0,1,6,7 / 0,1,3,4 & 3,3,9,9,3,3 \\
467 & 830|830|943|614|501|498 / 221200 & 432 / 8 & NO / 12 / B & 0,1,3,4,6,7 / 0,1,3,4,6,7 & 5,5,5,5,5,5 \\
468 & 614|501|498|830|830|943 / 200221 & 432 / 8 & ES / 12 / B & 0,1,3,4,6,7 / 0,1,3,4,6,7 & 5,5,5,5,5,5 \\
\end{longtable}
\endgroup

La tabla es una enumeración de la imagen de \(T_{\min}\). La especificación
del dominio permanece en \(\Omega_{\rm seed}=\mathcal S^2\); para recuperar
las condiciones iniciales individuales de una emisión se toma su fibra
\(T_{\min}^{-1}(U_6)\), cuya cardinalidad figura como
\(\#T_{\min}^{-1}(U_6)\).

\end{HMTScientificOwnerSubordinated}

\section{Certificación a profundidad \texorpdfstring{\(10^4\)}{10^4} de las expansiones coinductivas de \texorpdfstring{\(\pi\), \(e\) y \(\varphi\)}{pi, e y phi}}
\label{sec:truncamiento-10000-pi-e-phi-rev9}

\noindent\textbf{Alcance matemático.}
El teorema uniforme precedente cuantifica sobre todo entero positivo \(N\) y
construye el objeto infinito como límite inverso de sus cilindros compatibles.
Este capítulo ejecuta el caso \(N=10^4\) como certificado finito particularmente
exigente; \(10^4\) no es una cota del generador, del mismo modo que una ejecución
de un millón de cifras sería otra instancia del mismo algoritmo y no una
modificación de su ley.

La presente ejecución fija una profundidad de cálculo y comprueba que el
cilindro obtenido queda contenido en una única celda decimal de diez mil
posiciones. El menor número de bloques compatible con la monodromía nonádica
y con
\[
  3^{6B}>10^{10000}
\]
es
\[
  B=3501=389\cdot9,
  \qquad 6B=21006\ \text{trits}.
\]

Para cada uno de los tres canales se recorren exhaustivamente los \(729\)
subcilindros en cada etapa y se conserva el único subcilindro compatible con la horquilla
racional del lector. El cálculo produce
\[
 \begin{array}{c|r}
 \text{magnitud calculada por canal}&\text{valor}\\
 \hline
 \text{particiones completas }729\to1&3501\\
 \text{trits transportados}&21006\\
 \text{pares de igual fase }(K,K+9)&3492\\
 \text{vueltas nonádicas completas}&388
 \end{array}
\]
y no requiere una nonada suplementaria. Los primeros \(396\) bloques
coinciden exactamente con el testigo anterior de mil cifras.

\Needspace{28\baselineskip}
Los lectores de clausura, propagación y autoescala se evalúan mediante
desigualdades racionales exactas. La aceleración computacional reorganiza las
sumas por aritmética entera y no abre bibliotecas de constantes ni prefijos
decimales objetivo. Sólo después de escribir las salidas se ejecuta un
segundo cálculo independiente. Se comprueban, en este orden, la
generación racional de los tres canales, la composición G9--T1 de
\(\alpha_{\rm HMT}\) y la coincidencia de las diez mil cifras con las
expansiones de referencia consultadas sólo al final. Las tres comparaciones
son exactas en las diez mil posiciones calculadas. Esta comprobación finita
no sustituye el argumento uniforme de prolongación a profundidad arbitraria.

La prueba ampliada precisa también el significado del retorno. En los
\(3492\) pares por canal se conserva la fase y cambia siempre el estado
completo, pues avanzan profundidad, prefijo y cilindro. Una proyección local
reducida puede repetirse en algunos pares sin que retorne el estado TPK. Esta
distinción sustituye la afirmación más fuerte de aperiodicidad local por la
propiedad efectivamente demostrada: monodromía de fase con memoria del estado
enriquecido.

Las cuatro expansiones completas se imprimen en el
\cref{app:desarrollos-10000-rev9}. Las recurrencias, los intervalos
racionales y los conteos anteriores determinan el cálculo; las expansiones
impresas son sus truncamientos finitos y permiten contrastar cada cifra sin
convertirse en premisas del generador.

\chapter{Certificados de precisión arbitraria y muestras decimales finitas reproducibles}
\label{app:desarrollos-10000-rev9}

Este apéndice imprime las salidas completas de la ejecución descrita en
\cref{sec:truncamiento-10000-pi-e-phi-rev9,sec:alpha-t1-10000-rev9}. Los
saltos de línea son exclusivamente tipográficos. Cada expansión constituye
un testigo finito de la ley coinductiva de precisión arbitraria demostrada en
la Parte III; no define la constante ni limita la profundidad del generador.

\begingroup
\lstset{basicstyle=\ttfamily\tiny,breaklines=true,breakatwhitespace=false,
  frame=single,columns=fullflexible,keepspaces=true,numbers=none}

\section{Expansión coinductiva de \(\pi\) y certificado de truncación compatible}
\lstinputlisting{generated/decimals_10000_rev9/pi_10000_decimales_96col.txt}

\section{Expansión coinductiva de \(e\) y certificado de truncación compatible}
\lstinputlisting{generated/decimals_10000_rev9/e_10000_decimales_96col.txt}

\section{Expansión coinductiva de \(\varphi\) y certificado de truncación compatible}
\lstinputlisting{generated/decimals_10000_rev9/phi_10000_decimales_96col.txt}

\section{Expansión de \texorpdfstring{\(\alpha_{\rm HMT}\)}{alpha HMT}}
\lstinputlisting{generated/decimals_10000_rev9/alpha_T1_10000_decimales_96col.txt}
\endgroup


\chapter{Formalización asistida del núcleo algebraico}
\label{app:formalizacion-asistida-acumulativa}

La formalización comprueba las identidades finitas sobre matrices, álgebras y
estados dodecafásicos tipados. Su función es aislar las obligaciones
algebraicas susceptibles de reducción mecánica y mantener separadas la
procedencia de los datos, la enumeración finita y la prueba formal.

\HMTDemoteHeadings
\chapter{Verificación formal de identidades algebraicas finitas}
\label{app:lean}

Las igualdades finitas siguientes se formalizan en Lean y se deciden por
reducción del núcleo lógico, sin axiomas externos adicionales. Los teoremas
demuestran:
\begin{longtable}{p{.40\textwidth}p{.50\textwidth}}
\toprule
Teorema & Contenido\\
\midrule
\endhead
\texttt{aw\_symmetric}
& simetría de la matriz Paley--Witt \(A_W\);\\
\texttt{aw\_square\_minus\_identity}
& \(A_W^2=-I_6\) sobre \(\mathbb F_3\);\\
\texttt{x\_sq\_add\_one\_has\_no\_root}
& irreducibilidad de \(x^2+1\) sobre \(\mathbb F_3\);\\
\texttt{dual\_has\_nonzero\_nilpotent}
& existencia de un nilpotente en el álgebra parabólica;\\
\texttt{split\_has\_nontrivial\_idempotent}
& existencia de un idempotente no trivial en el álgebra hiperbólica;\\
\texttt{canonical\_charge}
& \(\sum_{m=1}^{12}K_m=6263\);\\
\texttt{hadamard\_orbit\_one/two/three}
& las tres órbitas de Hadamard producen el estado dodecafásico publicado.\\
\bottomrule
\end{longtable}

La formalización recibe las matrices y el vector dodecafásico como datos
tipados y verifica sus consecuencias algebraicas. La construcción del estado
firmado y la enumeración de (104\,976) semillas pertenecen a los teoremas
combinatorios previos y no se sustituyen por esta reducción formal.

\HMTRestoreHeadings

\chapter{Correspondencia entre enunciados, pruebas, certificados y propietarios científicos}
\label{app:indice-resultados-fuentes-acumulativo}

La tabla siguiente identifica la residencia matemática principal de cada
resultado y el control negativo que distingue su derivación de una mera
coincidencia. Ninguna fila reemplaza el enunciado ni la demostración principal.

\begingroup
\small
\setlength{\tabcolsep}{4pt}
\renewcommand{\arraystretch}{1.16}
\begin{longtable}{@{}>{\raggedright\arraybackslash}p{.20\textwidth}
  >{\raggedright\arraybackslash}p{.43\textwidth}
  >{\raggedright\arraybackslash}p{.27\textwidth}@{}}
\caption{Resultados, demostraciones de referencia y controles reproducibles.}
\label{tab:indice-fuentes-acumulativo}\\
\toprule
Resultado & Demostración de referencia & Control reproducible\\
\midrule
\endfirsthead
\toprule
Resultado & Demostración de referencia & Control reproducible\\
\midrule
\endhead
Construcción de APP, doble evaluación y salto espectral
& Capítulos de APP y aritmética genealógica
& Identidades de soporte, espectro y reconstrucción de las dos hojas\\
Estado TRIT, orientación y tres álgebras cuadráticas
& Capítulo de TRIT y clasificación algebraica
& Conservación de residuo, cociente, orientación y acarreo\\
Inventario operatorio completo del TPK
& Capítulos de transporte, memoria y operadores fundamentales
& Conmutatividad, naturalidad y prolongación compatible\\
Extrema y media razón holográfica
& Completación cúbica nonádica
& Identidad $1000=729+243+27+1$ y conservación de la unidad\\
Generación prolongable de \(\pi,e,\varphi\)
& Biografías coinductivas y monodromía nonádica
& Cilindros encajados, truncamientos compatibles y falsador de oráculo\\
Dos construcciones internas de \(\alpha_{\rm HMT}\)
& Rueda dodecafásica y cartas reversibles
& Equivalencia de salidas y separación del reconocimiento CODATA\\
Proyectores dodecafásicos y defecto \(-1/12\)
& Núcleo Hensel--Witt
& Simetría, idempotencia, rangos $3/4$ y traza normalizada\\
Doble proyección arquimediana y excepcional
& Lector de cilindros y lector de incidencia
& Mismo estado de entrada, codominios tipados y no identificación de salidas\\
Retículo de Leech, VOA y Moonshine
& Corredor excepcional completo
& Entrelazadores de incidencia y obstrucciones de equivarianza\\
Genealogía del número y funcionales espectrales
& Cierre del concepto de número y aritmética espectral
& Productos de Euler, caracteres y cotas dirigidas\\
Recta de acción y orden decimal \(-34\)
& Acción bidireccional y coordenadas angulares conjugadas
& Conúcleo de orden $16$ y separación del corrector relativo\\
Barbero--Immirzi, operador de área y espectro
& Tránsito hexada--octada y Hamiltoniano modular
& Canales $90/120$, área, primera ley y entropía relativa\\
Ángulos CKM y violación CP
& Capítulo de mezcla y reflexión racional
& Unitariedad, Jarlskog y cautela sobre la reflexión física\\
Electrón bidireccional y masa refinada
& Definición de partícula y primera realización electrónica
& Punto fijo central, lectores $K_D/K_\Omega$ y valor beta rector\\
Atlas \(81/56/13/324\) y operadores de masa
& Ley General de Masas y atlas ruta por ruta
& Cobertura completa, sector nulo y selección prospectiva\\
Permitividad y permeabilidad del vacío
& Respuesta constitutiva bicapa
& Reconstrucción separada de \(Z,\mu,\varepsilon,c\)\\
Formulación variacional y hamiltoniana de la gravedad
& Gravitación Einstein--Cartan--Holst
& Equivalencia Euler--Lagrange/Hamilton y controles torsionales\\
Navier--Stokes y Yang--Mills
& Clausuras terminales de las realizaciones fluida y de calibre
& Matriz PF84, coercividad y transferencia\\
\bottomrule
\end{longtable}
\endgroup

Las formulaciones reducidas anteriores no gobiernan los teoremas de la
Parte IV. El alcance de cada especialización se determina allí por su objeto,
su dominio, sus mapas y sus identidades de comparación. En particular, un
rótulo histórico no puede declarar por sí solo una conclusión abierta o
cerrada, del mismo modo que una clasificación documental no sustituye la
conmutatividad de un diagrama. Las construcciones finitas de la tabla anterior
se conservan como antecedentes; su prolongación se decide mediante los
comparadores explícitos de cada capítulo.

\chapter{Dominios, cartas e interfaces de realización}
\label{app:dominios-cartas-interfaces-acumulativo}

Este bloque separa los teoremas internos, las cartas de evaluación, las
interfaces físicas y los programas de extensión. La separación permite leer
cada igualdad en su dominio preciso y mantiene visible qué datos, calibres,
orientaciones o condiciones de frontera intervienen en una realización.

\HMTDemoteHeadings
\chapter{Dominios, codominios e hipótesis de los resultados}
\label{app:premisas}

\begin{enumerate}
\item El censo \(104\,976\to468\to243\), la imagen de 42 ternas, la
factorización \(468=26\times18\), las cinco regiones y los cinco ciclos son
teoremas exhaustivos del sistema finito de transporte especificado en esta
obra.

\item El catálogo enriquecido produce, por predicados finitos de fibra y sin
consulta decimal, la única órbita de clausura, la única órbita de
propagación y la única órbita de autoescala; el calibre interno orienta en
ellas los representantes \(010211,201101,121200\). Sus caracteres se
construyen, respectivamente, por la clausura pentarregional con retorno
unitario, la recurrencia semigrupal \((n+1)a_{n+1}=a_n\) y el rayo positivo
único de la incidencia \(F_{\rm av}\). La evaluación arquimediana publica
entonces los valores generados y el reconocimiento convencional los denomina
\(\pi,e,\varphi\). Para cada profundidad \(N\), la prolongación TPK produce
un prefijo único compatible con el de profundidad \(N-1\); el límite
inverso de esos prefijos fija la expansión completa. Por ello, las palabras,
los caracteres y los valores son etapas sucesivas de una construcción
interna, no etiquetas elegidas mediante tablas externas.

\item El teorema dodecafásico tiene por dominio causal el estado terminal
enriquecido \(x_{\rm term}\) producido por APP--TRIT--TPK. La carta firmada
publica desde él
\(x_{\rm term}\xrightarrow{R_{\rm sgn}}U_{12}^{\rm sgn}\), y la inversión
integral publica después
\(U_{12}^{\rm sgn}\xrightarrow{\mathcal H_{12}^{-1}}K\). En este dominio,
las correspondencias con \((D_3K,D_4K,Q)\), la representación de Hadamard y
la clausura dodecafásica son exactas y reversibles. Ni \(U\) ni \(K\) son
entradas generadoras. La traza reducida de una trayectoria posicional y el
sistema bivariado completo son objetos distintos, según se establece en el
\cref{chap:tpk-definicion}.

\item La igualdad con la cadena central 2018 es un reconocimiento externo.
La palabra matemática interna y su demostración son independientes de las
recomendaciones metrológicas posteriores.

\item La instancia finita de la doble proyección demuestra una bandera común,
\(W_{12}\) y \(M_{12}\). En la rama reticular se fijan el pegado de
\(A_2^{12}\), la cámara radial y el origen; el 3-vecino resultante se certifica
par, unimodular y sin raíces, y se identifica clásicamente con Leech. La
construcción de \(V^\natural\) y del Monstruo es una consecuencia clásica. La
flecha \(W_{12}\to W_{24}\) se realiza mediante el diseño residual de una
dodecada y la elevación única de sus hexadas, relativamente a una elección
\((D,\ell)\); las identificaciones se realizan mediante aplicaciones, no
mediante igualdades de cardinalidad. La compatibilidad proyectiva de estos
datos y su doble realización sobre el estado enriquecido común se demuestran
en los
\cref{p12-83-c17-s2:thm:lector-excepcional-global-rev6,p12-83-c17-s2:thm:doble-realizacion-global-rev6}.

\item El teorema de cilindros produce un único real para cada historia de su
dominio. La ley de prolongación cubre coherentemente cada compacto y el
agotamiento dirigido cubre la recta completa, como demuestran los
\cref{p12-83-c18-s1:thm:sobreyectividad-evaluacion-compacto,p12-83-c18-s1:prop:agotamiento-recta-real-u024,p12-83-c18-s2:hmtch:thm-cociente-real}.
Por tanto, la evaluación es sobreyectiva y su cociente arquimediano es
\(\mathbb R\); la cobertura es una conclusión del generador HMT, no una
premisa adicional.

\item El producto nativo está demostrado en el sector canónico de palabras.
La extensión monoidal al espacio completo de historias y el entrelazamiento
con las filtraciones de supervivencia se formulan como aplicaciones distintas
en el \cref{chap:producto-nativo}.
\end{enumerate}

\chapter{Invariancia causal de las interfaces respecto de referencias externas}
\label{chap:independencia-prospectiva}
\index{independencia prospectiva}\index{grafo causal}
\index{protocolo ciego}\index{torres!completitud reconstructiva}

\section{Criterio de independencia causal respecto de referencias externas: grafo dirigido acíclico y fijación previa de la salida}
\label{sec:definiciones-independencia-prospectiva}

Sean \(\mathsf P\) las primitivas internas, \(\mathsf R\) las reglas,
lectores, calibres y normalizaciones declarados, \(G\) un generador,
\(o=G(\mathsf P,\mathsf R)\) su salida y \(y\) una referencia externa.

\begin{definition}[Interfaz prospectivamente independiente]
\label{def:interfaz-prospectivamente-independiente}
La interfaz es prospectivamente independiente de \(y\) si:
\begin{enumerate}
 \item \(y\notin\operatorname{Anc}(o)\);
 \item \(\mathsf P,\mathsf R,G\) y los controles negativos se congelan antes
 de abrir \(y\);
 \item la ejecución no tiene acceso directo ni indirecto a \(y\);
 \item \(o\) y su traza se sellan antes de calcular una comparación
 \(\operatorname{cmp}(o,y)\).
\end{enumerate}
\end{definition}

\begin{definition}[Independencia de selección]
\label{def:independencia-seleccion}
La ceguera de cada ejecución no basta si, tras producir candidatos
\(G_j(\mathsf P,\mathsf R_j)\), se publica
\[
 j_\ast(y)=\arg\min_j
 \operatorname{cmp}\!\left(G_j(\mathsf P,\mathsf R_j),y\right).
\]
La independencia de selección prohíbe que esta operación determine la salida
publicada.
\end{definition}

\begin{definition}[Predicción multivaluada]
\label{def:prediccion-multivaluada}
Si los descriptores físicos no metrológicos y las restricciones internas
admiten varias multisecciones, la salida prospectiva es el conjunto completo
\[
 \widehat{\mathfrak R}(d)
 =\{\mathcal M:\mathcal M\text{ satisface todas las restricciones
 preinscritas para }d\}.
\]
La no unicidad es una salida del procedimiento. Elegir después un elemento
por proximidad a una magnitud está prohibido.
\end{definition}

\section{Teorema de no interferencia}
\label{sec:teorema-no-interferencia}

\begin{definition}[Grafo causal dirigido acíclico fiel y completo]
\label{def:dag-fiel-completo}
Un grafo dirigido acíclico \(D\) es fiel si cada nodo representa la misma
función y las mismas entradas que la construcción citada. Es completo para
una salida si toda dependencia causalmente pertinente aparece como arista o
raíz declarada. Fidelidad y completitud son hipótesis semánticas; la
aciclicidad del conjunto de datos no las demuestra.
\end{definition}

\begin{theorem}[Independencia prospectiva condicionada al grafo causal]
\label{thm:no-interferencia-prospectiva}
Sea un sistema determinista de ecuaciones estructurales representado por un
grafo dirigido acíclico finito \(D\), fiel y completo para \(o\). Si
\(y\notin\operatorname{Anc}_D(o)\), todas las raíces ancestrales de \(o\)
permanecen congeladas y la ejecución no tiene acceso oculto a \(y\), entonces
para cualesquiera \(y_0,y_1\),
\[
 o_{\operatorname{do}(y=y_0)}
 =o_{\operatorname{do}(y=y_1)}.
\]
\end{theorem}

\begin{proof}
Se ordenan topológicamente los antecesores de \(o\). Sus raíces son internas
y permanecen fijas. Por inducción sobre el orden, cada nodo recibe los mismos
argumentos bajo ambas intervenciones y, por determinismo, devuelve el mismo
valor. La igualdad se propaga hasta \(o\). El resumen criptográfico no crea
la independencia; sólo hace detectable una alteración posterior.
\end{proof}

\begin{corollary}[Permutación de la comparación externa]
\label{cor:permutacion-comparacion-externa}
Si se permutan nombres, filas o valores externos después del sellado, la
salida interna no cambia, aunque pueda cambiar su puntuación:
\[
 o_{\operatorname{do}(y=\sigma y_0)}
 =o_{\operatorname{do}(y=y_0)},\qquad
 \operatorname{cmp}(o,\sigma y_0)
 \ne\operatorname{cmp}(o,y_0).
\]
\end{corollary}

\begin{proposition}[Separación de ramas calibradas]
\label{prop:separacion-ramas-calibradas}
Si una referencia externa intervino en una selección histórica, la
dependencia no contamina los objetos internos que no descienden de esa
selección. La rama puede aislarse y someterse a un protocolo prospectivo
nuevo sin modificar el núcleo.
\end{proposition}

\begin{proof}
Ser antecesor es una propiedad local del grafo. Las ramas
\[
 Z_0\longrightarrow p=5\longrightarrow\text{escala electrónica},
\qquad
 \{\text{masas}\}\longrightarrow\{\text{firmas calibradas}\}
 \longrightarrow\text{tabla}
\]
quedan fuera del subgrafo protegido. Las torres, el carácter formal, el
atlas y el centro \((5,5)\) no reciben aristas desde esas referencias.
\end{proof}

\section{Invariancia contrafactual sobre el grafo causal dirigido y alcance de la conclusión}
\label{sec:certificado-independencia-prospectiva}

El grafo causal publicado contiene \(71\) nodos, \(89\) aristas,
\(16\) interfaces y
\(29\) salidas protegidas. La comprobación causal establece aciclicidad, recuentos,
rutas declaradas, existencia de las fuentes registradas, ausencia de
referencias externas entre los antecesores consignados e invariancia
contrafactual del modelo simbólico. El resultado es
\[
 \#\{\text{antecesores externos de salidas protegidas}\}=0.
\]
El certificado no demuestra por sí solo que no falte una arista ni audita
todas las lecturas de datos de cada programa científico. Su fuerza es, por
tanto, \textsc{demostrado con estructura de partida explícita}: la conclusión
se obtiene sobre el grafo dirigido acíclico expresamente declarado y queda
condicionada a su fidelidad y completitud.

La sustitución simultánea de las referencias CODATA, Planck,
Barbero--Immirzi, CKM, masas, gravedad y cristal temporal puede alterar las
comparaciones y las ramas calibradas. No altera, dentro del grafo declarado,
el catálogo de \(468\) emisiones, las cinco regiones de \(\pi\), la
generación coinductiva correlacionada, \(\alpha_{12}\), \(C^\ast\), la década \(-34\), las torres,
el atlas, el centro electrónico, la doble proyección ni el corredor
excepcional.

\section{Ensayo multivaluado sin acceso al valor objetivo y criterio de identificación}
\label{sec:protocolo-ciego-multivaluado}

\begin{algorithm}[Asignación futura sin magnitudes]
\label{alg:asignacion-futura-sin-magnitudes}
La interfaz física se ejecuta en cuatro fases irreversibles:
\begin{enumerate}
 \item \emph{Congelación}: se sellan autoridad, atlas, involuciones,
 restricciones, código, prioridad y controles negativos.
 \item \emph{Entrada permitida}: sólo se reciben sector, generación, carga,
 color, conjugación, espín/estadística, composición y tipo metrológico; se
 excluyen masas, firmas refinadas, residuos, ángulos y unidades.
 \item \emph{Salida}: se enumeran todas las multisecciones compatibles; para
 cada extensión de cada multisección se calculan ruta observable, registro,
 firma, \(\chi_{\rm form}\), especialización y observable. En objetos
 compuestos se empalman primero las extensiones y los registros. Se sella el
 multiconjunto completo.
 \item \emph{Apertura}: un custodio distinto abre las magnitudes y calcula
 residuos, incertidumbres y covarianzas sin modificar candidato alguno.
\end{enumerate}
\end{algorithm}

La multisección enriquecida, y no una extensión puntual seleccionada
retrospectivamente, es la salida del emparejamiento físico. La evaluación
posterior puede ejecutarse fibra a fibra. El algoritmo precedente es una
especificación falsable; no se presenta como una predicción multiconjunto ya
ejecutada.

\begin{criterion}[Fallo del protocolo ciego]
\label{crit:fallo-protocolo-ciego}
El protocolo falla si el código puede leer una magnitud, si una firma
histórica calibrada reaparece como entrada, si se descarta un candidato tras
abrir la referencia, si se confunde masa positiva con masa exactamente nula
o si el multiconjunto sellado cambia al permutar el catálogo externo.
\end{criterion}

\section{Carácter formal y completitud reconstructiva}
\label{sec:completitud-reconstructiva-torres}

Para \(v=(n_A,s_C,k,b_{120},b_{\zeta})\in\ZZ^5\), el carácter formal
anterior a toda especialización es
\[
 \chi_{\rm form}(v;x_A,x_C,x_\Delta,x_{120},x_{270})
 =\exp\!\left(
 n_Ax_A+s_Cx_C+kx_\Delta
 +b_{120}x_{120}+b_{\zeta}x_{270}\right).
\]
La especialización declarada
\[
 \Theta_0
 =\left(A_{\rm rad},\frac{C^\ast_{\rm rad}}6,
 \Dfour,s_{120},\zeta_{270}\right),
 \qquad \zeta_{270}=135\Delta_4^2,
\]
produce
\(\chi_0(v)
=\chi_{\rm form}(v;\Theta_0)\).
La firma cruda, el carácter formal, su especialización central, el
refinamiento global \(\eta\in\mathcal E_{\mathfrak S}\) y una tabla
calibrada son cinco objetos diferentes.

\begin{theorem}[Completitud reconstructiva positiva]
\label{thm:completitud-reconstructiva-positiva}
Sea
\[
 \mathcal N=\langle90,120\rangle_{>0}
\]
y sea \((n_j)_{j\ge1}\) su enumeración creciente. Dados
\(0<q_+<q_-<1\), definamos
\[
 a_j=q_-^{n_j}-q_+^{n_j}>0.
\]
Entonces \(a_j\to0\) y \(\sum_ja_j<\infty\). Fijada una regla determinista
para resolver empates de semienteros, para todo \(r\in\RR\) la recurrencia
\[
 r_0=r,\qquad
 \nu_j=\operatorname{nint}\!\left(\frac{r_{j-1}}{a_j}\right),\qquad
 r_j=r_{j-1}-\nu_ja_j
\]
satisface
\[
 |r_j|\leq\frac{a_j}{2}\longrightarrow0,\qquad
 r=\sum_{j\ge1}\nu_ja_j,\qquad
 \sum_{j\ge1}|\nu_ja_j|<\infty.
\]
Por tanto, todo \(R>0\), con \(r=\log R\), se representa exactamente en la
completación por
\(\exp(\sum_j\nu_ja_j)\).
\end{theorem}

\begin{proof}
Los índices de \(\mathcal N\) son \(90,120\) y todos los múltiplos de \(30\)
desde \(180\). En consecuencia,
\[
 \sum_{n\in\mathcal N}q_-^n
 =q_-^{90}+q_-^{120}
 +\frac{q_-^{180}}{1-q_-^{30}}<\infty.
\]
Como \(0<a_j<q_-^{n_j}\), la sucesión es sumable y converge a cero. La
definición de entero más próximo implica
\(|r_j|\le a_j/2\). Por telescopado,
\(\sum_{i=1}^{N}\nu_ia_i=r-r_N\to r\). Para \(j\ge2\),
\[
 |\nu_ja_j|=|r_{j-1}-r_j|
 \le\frac{a_{j-1}+a_j}{2},
\]
y la sumabilidad de \((a_j)\) prueba la convergencia absoluta.
\end{proof}

El teorema es una completitud del diccionario sobre \(\RR_{>0}\). Si los
coeficientes \(\nu_j\) se obtienen del residuo de una masa conocida, su uso
es \textsc{reconstrucción calibrada}; no proporciona la selección
prospectiva. Tampoco representa masa exactamente nula. Fotón, gluón o un
eventual modo de espín dos requieren protección exacta por simetría en la
representación física.

\Needspace{34\baselineskip}
\section{Protocolos prospectivos de invariancia para las familias de constantes}
\label{sec:protocolos-prospectivos-especificos}

\begin{description}
 \item[Barbero--Immirzi.] Se congelan \(S_{90},S_{120}\), los pesos
 \((12,-1)\) y su normalización antes de abrir una referencia física. La
 evaluación de \(\Gamma_{6\to8}\) es ciega respecto de esa referencia,
 pero la unicidad previa de series y pesos y el entrelazador con el operador
 de área son problemas distintos.
 \item[CKM y CP.] Se sellan las cuatro ecuaciones, sus coeficientes y la
 convención angular; se publican simultáneamente los cuatro parámetros,
 \(|V|\) y \(J\). El contraste usa una covarianza fijada de antemano.
 \item[Constitutivas.] Las identidades adimensionales preceden a
 \(Z_{\rm ref}\) y \(c_{\rm ref}\). El cambio de carta transforma la
 realización, no las identidades, y no crea cuatro predicciones
 independientes.
 \item[Gravedad e inercia.] Las razones
 \(G_{\rm tor}/I_{\rm tor}=1\) y
 \(G_{\rm av}/I_{\rm av}=\pi/\varphi^2\) pertenecen a capas diferentes.
 Su interpretación física exige un morfismo y un contraste prospectivo; un
 fenómeno cosmológico no selecciona la razón.
\end{description}

\section{Estratificación probatoria: identidades internas, deducciones, reconstrucciones calibradas e interfaces sometidas a cierre matemático}
\label{sec:cinco-niveles-independencia}

\begin{table}[htbp]
\centering
\caption{Niveles que la independencia causal no permite colapsar.}
\label{tab:cinco-niveles-independencia}
\small
\begin{tabularx}{\textwidth}{@{}L{0.25\textwidth}YL{0.27\textwidth}@{}}
\toprule
Nivel & Ejemplo & Afirmación permitida\\
\midrule
Generador interno certificado&
Regiones, \(\alpha_{12}\), \(C^\ast\), \(-34\), atlas&
Salida obtenida en el dominio declarado.\\
Evaluación condicionada ciega&
\(\Gamma_{6\to8}\), CKM fijado, \(\chi_0(v)\)&
No consulta el blanco, pero depende del funcional o firma.\\
Protocolo ciego no ejecutado&
Especificación multivaluada&
Existe un procedimiento falsable; no una salida sellada.\\
Reconstrucción calibrada&
Firmas históricas y refinamiento desde \(r=\log R\)&
Evaluación exacta condicionada a una referencia.\\
Interfaz física pendiente&
Partícula--antipartícula, área, masa nula, gravedad&
Dominio definido; falta representación o contraste.\\
\bottomrule
\end{tabularx}
\end{table}

\HMTRestoreHeadings

\chapter{Reproducibilidad del atlas y comparación de realizaciones de masa}
\label{app:reproducibilidad-masas-acumulativa}
\label{app:suplemento-digital-masas-rev11}

La teoría de masas se genera sobre el dominio completo de
\(81\) posiciones, \(56\) familias de ruta, \(13\) multisecciones,
\(324\) rutas orientadas y el sector nulo. La fase estructural calcula
\[
\text{posición o fibra}\longrightarrow\text{familia}
\longrightarrow\text{ruta orientada}\longrightarrow\text{registro}
\longrightarrow\text{firma nominal}\longrightarrow\text{carácter y torre},
\]
y, cuando la fibra posee el lector operatorial correspondiente, añade la
rama de refinamiento
\[
 \text{registro}\longrightarrow(K_D,K_\Omega)
 \longrightarrow\text{espectro de fibra}.
\]
La segunda rama refina la primera y no la sustituye: en particular,
\(W\), \(Z\) y \(H\) conservan sus rutas CT--108, sus firmas angulares y
sus caracteres escalares aunque no pertenezcan al atlas operatorial de
trece multisecciones. Estos invariantes se fijan antes del contraste externo;
una referencia metrológica no interviene en la selección de ruta, posición,
firma ni coeficiente.

\section{Invariantes finitos de la clasificación}
\label{sec:fuentes-certificado-tabla-masas-acumulativa}

\begin{itemize}
 \item atlas completo de $81$ celdas con sus lectores;
 \item $56$ familias de ruta con sus espectros;
 \item $13$ operadores de multisección;
 \item $324$ rutas orientadas con registro y carácter;
 \item inventarios completos de masas y anchuras;
 \item carta comparativa posterior a la selección genealógica.
\end{itemize}

\section{Condiciones de exactitud, estabilidad y reproducibilidad de los certificados}
\label{sec:ejecucion-tabla-masas-acumulativa}

El control integral exige la cobertura \(81/56/13/324\), el sector nulo, los
\(14\) perfiles de \(K_D\), los \(9\) perfiles de \(K_\Omega\), los
\(40\) pares conjuntos y las \(18\) coincidencias de lectores. Comprueba
además que la masa electrónica refinada se evalúa desde una única expresión,
que los sectores nulos conservan \(\kappa\) no aplicable, que una fila
operatoria no publica un residual escalar, que las doce rutas nominales
\(e,\mu,\tau,u,d,s,c,b,t,W,Z,H\) conservan su carácter escalar y que la
realización física mantiene el codominio propio de cada salida. La falta de un par
\(K_D/K_\Omega\) en un sector no puede degradar ni borrar su lector escalar
nativo.

La salida electrónica reproducida es
\[
m_e^\beta
=m_e^{(0)}R_{\rm act}^{80-54\alpha_{\rm HMT}+6\Delta_4}
=0.510998950690000808608257165\ldots\;\mathrm{MeV}/c^2.
\]
Las fibras de rango mayor publican su par operatorio, sus espectros y el
morfismo de realización correspondiente. La selección de autovalores queda
determinada por ese morfismo antes de la comparación externa.

\enlargethispage{6\baselineskip}
\section{Dominio científico de la clasificación}
\label{sec:alcance-publico-apendices-acumulativo}

El dominio de masas es el atlas completo anterior. La construcción recibe
exclusivamente posiciones, familias, multisecciones, rutas, registros y
operadores tipados; los datos de procedencia permanecen separados de la carga
científica.
